\documentclass[output=paper,colorlinks,citecolor=brown]{langscibook}
\ChapterDOI{10.5281/zenodo.21237007}
\author{Miguel Villanueva Svensson\orcid{}\affiliation{Vilnius University}}
%\ORCIDs{}

\title{Balto-Slavic denominatives from an Indo-European perspective: An update}

\abstract{This paper presents an updated and critical overview of the Balto-Slavic denominative formations from an Indo-European perspective. Some important points: 1) denominatives in *-\textit{ō-i̯e/o}- and *-\textit{ā-i̯e/o}- were still directly linked to \textit{o}- and \textit{ā}-stem nouns, respectively; 2) Balto-Slavic has no “\textit{ā}-factitives”; 3) we have a far better understanding of the development of PIE *-\textit{ei̯e/o}-denominatives (viz. causative-iteratives) in Balto-Slavic than just a few years ago; 4) we have new accounts of some suffixes; 5) Balto-Slavic deadjectives included not only essives and fientives, but also specialized factitives (most likely continuing PIE \textit{nu}-factitives); 6) after the breakup of Balto-Slavic, Baltic and Slavic evolved in different directions, in part due to different areal stimuli; 7) issues of directionality in the interplay between denominatives and deverbatives are also discussed.}


\IfFileExists{../localcommands.tex}{
   \addbibresource{../localbibliography.bib}
   \usepackage{tabularx,multicol}
\usepackage{url}
\urlstyle{same}

 
\usepackage{langsci-optional}
\usepackage{langsci-lgr}
\usepackage{langsci-gb4e}
\usepackage[linguistics,edges]{forest}
\usepackage{tikz-qtree}
\usetikzlibrary{positioning}
\usetikzlibrary{patterns.meta}
\usepackage{multirow}
\usepackage{pgfplots}
\usepackage{setspace}
\usepackage{amsmath,stackengine}
% \usepackage{fontspec}
% \usepackage{tipa} % TIPA is banned
\usepackage{langsci-textipa}
\usepackage{langsci-branding}
\usepackage{longtable}
\usepackage{tabto}

\usepackage{enumitem}

%\usepackage{lscape} %Querformat chapter 4

   % \newcommand*{\orcid}{}


\makeatletter
\let\thetitle\@title
\let\theauthor\@author
\makeatother

\newcommand{\togglepaper}[1][0]{
   \bibliography{../localbibliography}
   \papernote{\scriptsize\normalfont
     \theauthor.
     \titleTemp.
     To appear in:
    Laura Grestenberger, Viktoria Reiter \& Melanie Malzahn (eds.).
    \emph{Deadjectival verb formation in Indo-European and beyond}.
     Berlin: Language Science Press. [preliminary page numbering]
   }
   \pagenumbering{roman}
   \setcounter{chapter}{#1}
   \addtocounter{chapter}{-1}
}

\newbool{bookcompile}
\booltrue{bookcompile}
\newcommand{\bookorchapter}[2]{\ifbool{bookcompile}{#1}{#2}}

%Sonderzeichen
%idg palatales k
\newcommand{\kinvbreve}{k\raisebox{0.6ex}{\hspace{-0.05em}\char"0311}}

\newcommand{\jac}{%
  \textit{\char"0237}% italic dotless j (ȷ)
  \hspace{-0.01em}% fine-tune horizontal spacing
  \raisebox{0.01ex}{\char"0301}% raised combining acute
}

%Baltic circumflex l
\newcommand\ltilded{%
\hspace{-0.2em}%
  \stackon[-1.5pt]{l}{\smash{\kern0.2em\~{}}}%
  \hspace{-0.2em}%
}

%Slav. wedge+ double acute
\newcommand{\HighH}[1]{%
    \stackon[-1.2ex]{#1}{\kern0.1em\H{}\kern-0.1em}%
}

%Lith dot tilde VICKY diese b.-sl. Sonderzeichen BRINGEN MICH UM!!
\newcommand{\tildedot}[1]{%
    \ensurestackMath{%
        \stackon[-0.27ex]{% Tilde layer 
            \stackon[-0.15ex]{% Dot layer 
                #1%
            }{\hspace{0.1em}\cdot\hspace{-0.13em}}%
        }{\hspace{0.1em}\text{\~{}}\hspace{-0.13em}%
    }%
}}

%Lith dot acute
\newcommand{\dotacute}[1]{%
    \ensurestackMath{%
        \stackon[-1.1ex]{% acute layer 
            \stackon[-0.15ex]{% Dot layer 
                #1%
            }{\hspace{0.1em}\cdot\hspace{-0.12em}}%
        }{\hspace{0.1em}\text{\'{}}\hspace{-0.12em}%
    }%
}}



% \setlength{\emergencystretch}{2pt} 

 \pgfplotsset{compat=1.18}


\defcitealias{AbB_1}{AbB 1}
\defcitealias{AbB_7}{AbB 7}
\defcitealias{ABL_1}{ABL 1}
\defcitealias{Vries1977}{AEW}
\defcitealias{ALEW}{ALEW}
\defcitealias{AnGr}{AnGr}
\defcitealias{ARM_1}{ARM 1}
\defcitealias{BT}{BT}
\defcitealias{black_concise_2000}{CDA}
\defcitealias{CHD}{CHD}
\defcitealias{CT28}{CT 28}
\defcitealias{CT34}{CT 34}
\defcitealias{CT39}{CT 39}
\defcitealias{Cleasby1874}{CV}
\defcitealias{Beekes2010}{EDG}
\defcitealias{eDiAna}{eDiAna}
\defcitealias{eDIL}{eDIL}
\defcitealias{EDL}{EDL}
\defcitealias{Kroonen2013}{EDPG}
\defcitealias{ESJS}{ESJS}
\defcitealias{EVP}{EVP}
\defcitealias{EWA1988}{EWA}
\defcitealias{EWAI}{EWAia I}
\defcitealias{EWAII}{EWAia II}
\defcitealias{EWGP}{EWGP}
\defcitealias{GDLI1961}{GDLI}
\defcitealias{GPC}{GPC}
\defcitealias{GRADIT1999}{GRADIT}
\defcitealias{HW}{HW}
\defcitealias{IEW}{IEW}
\defcitealias{KAH_2}{KAH 2}
\defcitealias{KAR_1}{KAR 1}
\defcitealias{KAR_2}{KAR 2}
\defcitealias{KEWA}{KEWA}
\defcitealias{Labat_TDP}{Labat TDP}
\defcitealias{lambert_BWL}{Lambert BWL}
\defcitealias{LEIA}{LEIA}
\defcitealias{Rix2001}{LIV\textsuperscript{2}}
\defcitealias{LIVaddenda}{LIV\textsuperscript{2} {A}ddenda}
\defcitealias{LKG1971}{LKG}
\defcitealias{LKZ}{LKŽ}
\defcitealias{LP}{LP}
\defcitealias{MhdGr}{MhdGr}
\defcitealias{MLLVG1959}{MLLVG}
\defcitealias{NIL}{NIL}
\defcitealias{OgnS}{OgnS}
\defcitealias{ONP}{ONP}
\defcitealias{RIMA_2}{RIMA 2}
\defcitealias{TIM_2}{TIM 2}
\defcitealias{WAP}{WAP}
\defcitealias{YOS_2}{YOS 2}






\newindex{wan}{wdx}{wnd}{Form index}
\newcommand{\iw}[2]{\index[wan]{#1!{#2}\xspace}}
\newcommand{\iwi}[3][]{\index[wan]{#2!{#3}\xspace#1}{#3}}
\newcommand{\iwit}[3][]{\rule{0pt}{0pt}#1\index[wan]{#2!{#3}}{#3}}


\newindex{van}{vdx}{vnd}{Form index}
\newindex{ban}{bdx}{bnd}{Book index}

% \newcommand{\iw}[3][]{\index[wan]{#2!\textit{#3}}}


% \providecommand{\iwi}[3][]{\iw{#3}\textit{#3}}


\newcommand{\indexnote}[1]{{\footnotesize\sffamily\raggedright\normalfont\color{gray}#1}}

   %% hyphenation points for line breaks
%% Normally, automatic hyphenation in LaTeX is very good
%% If a word is mis-hyphenated, add it to this file
%%
%% add information to TeX file before \begin{document} with:
%% %% hyphenation points for line breaks
%% Normally, automatic hyphenation in LaTeX is very good
%% If a word is mis-hyphenated, add it to this file
%%
%% add information to TeX file before \begin{document} with:
%% %% hyphenation points for line breaks
%% Normally, automatic hyphenation in LaTeX is very good
%% If a word is mis-hyphenated, add it to this file
%%
%% add information to TeX file before \begin{document} with:
%% \include{localhyphenation}
\hyphenation{
    par-a-digm non-de-adj-ect-iv-al -pre-sent -pre-sents for-ma-tion-al in-do-ger-ma-ni-sche in-do-ger-ma-ni-schen Rei-chert Rei-ter dia-chrony ety-mo-lo-gi-sches Ost-ger-ma-ni-schen ger-ma-ni-schen alt-in-di-schen mor-pho-lo-gi-scher Alt-indo-ari-sch-en lexico-graphy Schrij-ver
}
% aktuell in Bibliographie: stud-ies, histor-ical Soci-età
\hyphenation{
    par-a-digm non-de-adj-ect-iv-al -pre-sent -pre-sents for-ma-tion-al in-do-ger-ma-ni-sche in-do-ger-ma-ni-schen Rei-chert Rei-ter dia-chrony ety-mo-lo-gi-sches Ost-ger-ma-ni-schen ger-ma-ni-schen alt-in-di-schen mor-pho-lo-gi-scher Alt-indo-ari-sch-en lexico-graphy Schrij-ver
}
% aktuell in Bibliographie: stud-ies, histor-ical Soci-età
\hyphenation{
    par-a-digm non-de-adj-ect-iv-al -pre-sent -pre-sents for-ma-tion-al in-do-ger-ma-ni-sche in-do-ger-ma-ni-schen Rei-chert Rei-ter dia-chrony ety-mo-lo-gi-sches Ost-ger-ma-ni-schen ger-ma-ni-schen alt-in-di-schen mor-pho-lo-gi-scher Alt-indo-ari-sch-en lexico-graphy Schrij-ver
}
% aktuell in Bibliographie: stud-ies, histor-ical Soci-età
   \boolfalse{bookcompile}
   \togglepaper[5]%%chapternumber
}{}


\begin{document}
\SetupAffiliations{mark style=none}
\maketitle

\section{Introduction}
\label{sec:BSlIntro}

The purpose of this article is to present an overview of Balto-Slavic denominatives\is{denominative!verb} from an Indo-European perspective, stressing recent findings and pending issues. It should be noted from the outset that the focus is on the Balto-Slavic system, not on PIE reconstruction or on the historical Baltic and Slavic languages (though these are, needless to say, treated as well). The article is divided into two parts: survey of Balto-Slavic denominative formations\is{denominative!verb} (Section \ref{sec:BSldenom}); discussion of their relevance for Indo-European studies (Section \ref{sec:BSlissues}).

Before proceeding further, a note on terminology. \textit{Deverbative}\is{deverbal/deverbative!verb} designates a verb derived from a verb, not a \textit{primary} \textit{verb}.\is{derivation!deradical} \textit{Denominative}\is{denominative!verb} is used as a cover term for both \textit{desubstantive} and \textit{deadjective}.\is{derivation!deadjectival} The latter is regularly so specified; the former only occasionally. I distinguish between \textit{stative},\is{stative} \textit{inchoative}\is{inchoative} and \textit{causative}\is{causative} for deverbatives\is{deverbal/deverbative!verb} and \textit{essive},\is{essive} \is{fientive} \textit{fientive}\footnote{The terms “\isi{essive}” and “\isi{fientive}” as used here should not be confused with their use for the alleged suffixes “\iwi{PIE}{*-\textit{éh\textsubscript{1}-/-h\textsubscript{1}}-}”, “\iwi{PIE}{*-\textit{h\textsubscript{1}-i̯e/o}-}” in part of the Indo-European literature (e.g., \citetalias[ 25]{Rix2001}).} and \textit{factitive}\is{factitive} for denominatives.\is{denominative!verb} \textit{Iterative}\is{iterative} is sometimes used as a cover term for several semantically differentiated deverbative\is{deverbal/deverbative!verb} subtypes. As for language labels, most of them require no comment. I use \textit{Core Indo-European} for the parent language after the separation of Anatolian, \textit{North-Western Indo-European} for the area of Indo-European from which Italo-Celtic, Germanic and Balto-Slavic derive, and \textit{Northern Indo-European}\is{Northern Indo-European} for the common ancestor of Germanic and Balto-Slavic. None of these entities is uncontroversial, but a detailed treatment would be out of place here. \textit{Balto-Slavic} unity is simply taken for granted.

The term “update” in the title should be understood in very vague terms. Denominatives\is{denominative!verb} are not a popular research topic and are often absent from short surveys.\footnote{Hardly anything could be grasped from \citet{Hock2005} or the relevant chapters of \citet{Kapovic2017} or \citet{Klein2018}.} Handbooks like \citet[86--93]{Birnbaum1997} or \citet[\textit{passim}]{Schmalstieg2000} are more informative, but are rarely up to date from an Indo-European perspective. Thus, there is no real \textit{terminus a quo} for the update, though as a rule I do not quote comparative literature published before the classical grammars of \citet{Stang1966} and \citet{Vaillant1966}. Finally, any “update” of this sort is per force going to be personal and mine is no exception. I have nevertheless made a conscious effort not to miss alternative views.

\section{Balto-Slavic denominative formations}
\label{sec:BSldenom}

\subsection{Introduction}
\label{sec:BSldenomintro}

I first present desubstantives (\ref{sec:BSlaje}-\ref{sec:BSli}), then deadjectives \is{deadjectival!verb} (\ref{sec:BSleje}-\ref{sec:BSlneu}). Needless to say, a sharp distinction between the two is impossible. Ideally, each entry includes the following information: 

\begin{enumerate}
    \item Balto-Slavic suffix; 
    \item Baltic and Slavic evidence;\footnote{I regularly give the infinitive, present and aorist/preterit. Lithuanian inflected forms are given in the 3\textsuperscript{rd} person, Latvian and OCS forms in the 1\textsuperscript{st} singular (or, in the case of OCS, sometimes just the stem). Quotation of Old Prussian forms naturally depends on what is actually attested.}
    \item comments on the Balto-Slavic prototype;
    \item PIE background;
    \item Balto-Slavic contribution.
\end{enumerate}

For considerations of space, quotation of evidence is reduced to a minimum. The same holds true for references to the literature. In particular, I usually omit references to the standard grammars and other reference works.\footnote{Descriptive surveys of the East Baltic languages are not hard to find. The best ones for the present topic probably remain the academic grammars produced in the mid-20\textsuperscript{th} century (\citetalias[ 330--343]{MLLVG1959}; \citetalias[ 247--268]{LKG1971}). My own presentation of the facts is mostly based on these sources (and, for Lithuanian, on \citealt{Jakaitiene1973}). I regret to say that I have found more recent grammars far less informative. On the other hand, older sources like \citet[488--551]{Skardzius1943} and, especially, \citet[618--651]{Endzelin1923} remain invaluable. For Old Prussian, I mostly rely on \citet[107--114]{Endzelin1943}. See also \citet{Smoczynski2005} for the data. \citet[360--374]{Stang1966} remains unsurpassed as a \textit{global} comparative treatment. The case of Slavic is different. For comparative work, only the very securely reconstructed Common Slavic matters, which can often be informally represented with attested Old Church Slavonic. \citet{Vaillant1966} not only remains unsurpassed; all later general treatments I have consulted are less informative. As a result, my presentation of the Slavic facts closely follows that of Vaillant.} Note that while all three Baltic languages are represented, Slavic is usually limited to Old Church Slavonic as an informal representation of Late Common Slavic. Common Slavic reconstructions are only provided when Old Church Slavonic is uninformative (most saliently, for accentology).\footnote{Acuteness is marked as *\underline{Ē}, with underline, and non-acute as simple *Ē. Accentology is otherwise not treated in this survey.} The modern Slavic languages stand outside the scope of this paper. Finally, note that for Balto-Slavic I distinguish between (e.g.) “\textit{ī}-presents” and “\textit{ī}-verbs”. This is unnecessary for Indo-European, where suffixes like *-\textit{ei̯e/o}-\iw{PIE}{*-\textit{éi̯e/o}-} were limited to the present stem.

\subsection{Bl.-Sl.\ *-\textit{\underline{ā}-i̯e/a}- (inf.\ *-\textit{\underline{ā}-tī}, aor.\ *-\textit{\underline{ā}-s}- [\textit{vel sim}.])}
\label{sec:BSlaje}

This suffix is well represented in the three Baltic languages: Lith.\ \textit{núom-oti}, -\textit{oja}, \textit{-ojo}\iw{Lith}{\textit{núomoti, -oja, -ojo}} ‘rent’ ($\leftarrow$ \iwi{Lith}{\textit{núoma}} ‘rent’); Latv.\ \textit{nuõm-ât}, -\textit{ãju}, -\textit{ãju}\iw{Latv}{\textit{nuõmât, -ãju, -ãju}} ‘id.’ ($\leftarrow$ \iwi{Latv}{\textit{nuõma}} ‘id.’); OPr.\ \textit{dwigub-ūt}, \textit{dwigubb-ū}\iw{OPr}{\textit{dwigubūt, dwigubbū}}  ‘doubt’ ($\leftarrow$ fem. gen. sg.\ \textit{dwigubbus}\iw{OPr}{\textit{dwigubbus} (gen.sg.f.)} ‘twofold’). It regularly builds verbs from nouns, only occasionally from adjectives. The derivational dependency from \textit{ā}-stems is still well preserved in Lithuanian (71,53\% of the denominatives\is{denominative!verb} in \iwi{Lith}{\textit{-oti}} are derived from \textit{ā}-stem nouns, according to \citealt[281]{Pakerys2011}). I do not have exact numbers for Latvian, but the forms given in the grammars point in the same direction. This suffix is also productive\is{productivity} for deverbative\is{deverbal/deverbative!verb} \is{iterative}iteratives: Lith.\ \textit{šaũkti}, -\textit{ia}\iw{Lith}{\textit{šaũkti, -ia}} ‘shout’ $\rightarrow$ \textit{šaũk-oti}, -\textit{oja}\iw{Lith}{\textit{šaũkoti, -oja}} ‘id. (iter.)’, Latv.\ \iwi{Latv}{\textit{sàukt}} $\rightarrow$ \iwi{Latv}{\textit{saũkât}}. According to \citet[34]{Jakaitiene1973}, in Lithuanian, deverbatives\is{deverbal/deverbative!verb} are twice as common as the denominatives,\is{denominative!verb} though usually in extended forms like \iwi{Lith}{-\textit{čioti}}, \iwi{Lith}{-\textit{soti}}, \iwi{Lith}{-\textit{n(i)oti}}, \iwi{Lith}{-\textit{l(i)oti}} or, above all, \iwi{Lith}{-\textit{ioti}} (\textit{b\dotacute{e}gti}, -\textit{a}\iw{Lith}{\textit{b\dotacute{e}gti, -a}} ‘run’ $\rightarrow$ \textit{b\.egióti}, -\textit{iója}\iw{Lith}{\textit{b\.egióti, -iója}} ‘id. (iter.)’; see \ref{sec:BSli}, \ref{sec:BSlafter}).

Slavic essentially concords with Baltic, with the important difference that \iwi{OCS}{\nobreakdash-\textit{ati}, -\textit{ajǫ}} is immensely productive\is{productivity} for derived \isi{imperfective}s (OCS \iwi{OCS}{\textit{čisti}, \textit{čьtǫ}} ‘read’ $\rightarrow$ \textit{po-čitati}, -\textit{čitajǫ}\iw{OCS}{\textit{počitati, -čitajǫ}} ‘id. (impf.)’). As a result, denominatives\is{denominative!verb} are very well represented, but are not productive\is{productivity} anymore in most Slavic languages: OCS \iwi{OCS}{\textit{igra}} ‘play’ $\rightarrow$ \textit{igr\nobreakdash-ati}, -\textit{ajǫ}\iw{OCS}{\textit{igrati, -ajǫ}} ‘play’, \iwi{OCS}{\textit{dělo}} ‘work’ $\rightarrow$ \textit{děl-ati}, -\textit{ajǫ}\iw{OCS}{\textit{dělati, -ajǫ}} ‘do’. I do not have numbers indicating whether the connection with \textit{ā}-nouns is as clear as in Baltic. Deadjectives occur, but are very rare.\is{deadjectival!verb}

For Balto-Slavic we can safely reconstruct a denominative suffix\is{denominative!suffix} *-\textit{\underline{ā}-tī}, *-\textit{\underline{ā}\nobreakdash-i̯e/a}-\iw{PBSl}{*-\textit{ā-tī}, *-\textit{ā-i̯e/a}-} that was still derivationally attached to \textit{ā}-stem nouns. It is impossible to say how productive\is{productivity} derivation from other stems was. It was probably not used for deadjectives,\is{deadjectival!verb} but deverbative\is{deverbal/deverbative!verb} iteratives\is{iterative} were very productive.\is{productivity}

The origin of this type is perfectly clear: It goes back to PIE \textit{i̯e/o}-denominatives\is{denominative!verb} \is{*-\textit{i̯e/o}-present} from \textit{ah\textsubscript{2}}-stem nouns. Whether *-\textit{ah\textsubscript{2}-i̯e/o}-denominatives\is{denominative!verb} expanded beyond their original core within PIE is a matter I will leave open.\footnote{See \citet[126--127]{Sasseville2021} with references for recent discussion.} In any case, it is interesting to note that the original connection with the \textit{ā}-stems is still apparent in Lithuanian. Another observation to make is that, although iteratives\is{iterative} in \iwi{PIE}{*-\textit{ah\textsubscript{2}-i̯e/o}-} are found in several languages and may even go back to the parent language, they are nowhere as productive\is{productivity} as in Germanic and Balto-Slavic.

Before moving on, a note on the \emph{ā-presents}. As is well known, Baltic is one of the few languages that has an \iwi{PB}{*-\textit{ā}-}present distinct from \iwi{PIE}{*-\textit{ā-i̯e/o}-} (Lith.\ \textit{laĩko}\iw{Lith}{\textit{laikýti, laĩko, laĩk\.e}, 1sg. prs. \textit{laikaũ}} ‘holds’ vs. \iwi{Lith}{\textit{galvója}} ‘thinks’). The other branches in which this class is attested are Anatolian (Hitt.\ \textit{šuppiyaḫḫ-\textsuperscript{ḫḫi}}  \iw{Hitt}{\textit{šuppiyaḫḫ}-} ‘purify’; cf.\ \citealt[140--141]{Jasanoff2003}) and Greek (cf.\ \citealt{Rau2009myc}). This type has recently been put in a new light by \citegen{Sasseville2015} demonstration that in Luvian *-\textit{ah\textsubscript{2}}-verbs built denominatives\is{denominative!verb} from *-\textit{ah\textsubscript{2}}-nouns through \isi{conversion}. It follows that the now standard assumption of a class of “\textit{ah\textsubscript{2}}-\isi{factitive}s of the type \textit{renouāre}\iw{Latin}{\textit{(re)nouāre}} derived from \textit{o}-stems” (\textit{vel sim}.) is in need of revision. Here I will limit myself to stress that, against what is often assumed, \textit{ā}-presents are exclusively deverbative\is{deverbal/deverbative!verb} in Balto-Slavic; cf.\ \citet[392--395]{VillanuevaSvensson2020}.

\subsection{Bl.-Sl.\ *-\textit{\underline{ō}-i̯e/a}- (inf.\ *-\textit{\underline{ō}-tī}, aor.\ *-\textit{\underline{ō}-s}- [\textit{vel sim}.])}
\label{sec:BSloje}

This suffix is only unambiguously attested in East Baltic, where it is very productive:\is{productivity} Lith.\ \textit{mel-úoti}, -\textit{úoja}, -\textit{ãvo}\iw{Lith}{\textit{melúoti, -úoja, -ãvo}} ‘lie’ ($\leftarrow$ \iwi{Lith}{\textit{mẽlas}} ‘lie’), Latv.\ \textit{mel-uôt}, -\textit{uõju}, -\textit{uõju}\iw{Latv}{\textit{meluôt, -uõju, -uõju}} ‘id.’ ($\leftarrow$ \iwi{Latv}{\textit{mȩli}} [pl.t.]). In Lithuanian, verbs in \iwi{Lith}{-\textit{úoti}, -\textit{úoja}} are normally made from \textit{o}-stem nouns (70,95\% of the denominatives,\is{denominative!verb} according to \citealt[286]{Pakerys2011}). I do not have numbers for Latvian, but the impression one gets from the grammars is that the connection with \textit{o}-stem nouns is not as clear as in Lithuanian. Deadjectives\is{deadjectival!verb} are rarer, but not exceptional: Lith.\ \iwi{Lith}{\textit{báltas}} ‘white’ $\rightarrow$ \textit{balt-úoti}, -\textit{úoja}\iw{Lith}{\textit{baltúoti, -úoja}} ‘show white’, Latv.\ \iwi{Latv}{\textit{skaĩdrs}} ‘clear’ $\rightarrow$ \textit{skaĩdr-uôt}, -\textit{uõju}\iw{Latv}{\textit{skaĩdruôt, -uõju}} ‘make clear’, refl. ‘become clear’. Their meaning is broadly the same as that of the deadjectives \is{deadjectival!verb} in *-\textit{\underline{ē}ti}\iw{PBSl}{*-\textit{ē-tī}, *-\textit{ē-i̯e/a}-} (\isi{fientive} or \isi{essive}, occasionally also \isi{factitive}). Deverbatives\is{deverbal/deverbative!verb} are limited to Lithuanian and clearly secondary\is{derivation!stem-based} (\textit{a\ltilded pti}, -\textit{sta}\iw{Lith}{\textit{a\ltilded pti, -sta}} ‘faint’ $\rightarrow$ \textit{alp-úoti}, -\textit{úoja}\iw{Lith}{\textit{alpúoti, -úoja}} ‘id. (iter.)’).

Due to the merger of Bl.-Sl.\ *\textit{ō} and *\textit{ā} in West Baltic and Slavic, it is impossible to know whether these branches had verbs in *-\textit{\underline{ō}-tī}, *-\textit{\underline{ō}-i̯e/a}-\iw{PBSl}{*-\textit{ō-tī}, *-\textit{ō-i̯e/a}-} beside the familiar *-\textit{\underline{ā}-tī}, *-\textit{\underline{ā}-i̯e/a}-\iw{PBSl}{*-\textit{ā-tī}, *-\textit{ā-i̯e/a}-}. If *-\textit{\underline{ō}-tī}, *-\textit{\underline{ō}-i̯e/a}-\iw{PBSl}{*-\textit{ō-tī}, *-\textit{ō-i̯e/a}-} is projected into Balto-Slavic, it certainly built denominatives\is{denominative!verb} from \textit{o}-stem nouns. It is impossible to know whether it also built desubstantives from other stems and deadjectives.\is{deadjectival!verb}

The only potential \textit{comparandum} of Bl.-Sl.\ *-\textit{\underline{ō}-tī}, *-\textit{\underline{ō}-i̯e/a}-\iw{PBSl}{*-\textit{ō-tī}, *-\textit{ō-i̯e/a}-} is the Greek type \iwi{AG}{δουλόω} ‘enslave’ ($\leftarrow$ \iwi{AG}{δοῦλος} ‘slave’).\footnote{This statement obviously implies that the short stem vowel of Gk.\ δουλ-ό-ω\iw{AG}{δουλόω} (and of \iwi{AG}{τῑμάω} ‘honor’ or \iwi{AG}{ἀνθέω} ‘be in bloom’) is an inner-Greek development (contrast aor.\ \iwi{AG}{ἐδούλωσα}, \iwi{AG}{ἐτῑ́μησα}, \iwi{AG}{ἤνθησα}), as is the \textit{communis opinio}. The issue cannot be discussed at length here.} There is, however, no \textit{communis opinio} concerning the connection between the two formations (if any). During the last decades several scholars have proposed reconstructing for PIE a denominative suffix\is{denominative!suffix} \iwi{PIE}{*-\textit{oh\textsubscript{1}-i̯e/o}-} (\citealt{Peters1999}; \citealt[401--402]{Malzahn2010}) or \iwi{PIE}{*-\textit{o-i̯e/o}-} (\citealt[357--358]{Oettinger1979}; \citealt[132--133]{Kloekhorst2008}). The evidence supporting such a view, however, is either inherently problematic (e.g., Gaulish \iwi{Gaulish}{\textit{marcosior}} ‘I will ride/I will be ridden’) or is better explained in a different way (the Hittite \iwi{Hitt}{\textit{ḫatrā(i)}-}type, for instance, is usually derived from \iwi{PIE}{*-\textit{ah\textsubscript{2}-i̯e/o}-}; see \citealt[116--123]{Sasseville2021}, with references). The dominant position has always been that Lith.\ \iwi{Lith}{-\textit{úoti}, -\textit{úoja}} and Gk.\ \iwi{AG}{-όω} are parallel, but independent innovations; see especially \citet[15--22]{Tucker1981} for the latter. Within this approach Lith.\ \iwi{Lith}{-\textit{úoti}, -\textit{úoja}} is usually derived from \isi{possessive} adjectives in “\iwi{PIE}{*-\textit{ō-to}-}” of the type Gk.\ \iwi{AG}{χολωτός} ‘angry’, Lith.\ \iwi{Lith}{\textit{ragúotas}} ‘horned’ (e.g., \citealt[627--628]{Endzelin1923}; \citealt[364]{Stang1966}) and this remains the best option today.

I add a few observations on Balto-Slavic:

\begin{itemize}
    \item it is impossible to know the chronology of the innovation and whether *-\textit{\underline{ō}-tī}, *-\textit{\underline{ō}-i̯e/a}-\iw{PBSl}{*-\textit{ō-tī}, *-\textit{ō-i̯e/a}-} is really Balto-Slavic in date;
    \item the now generally accepted derivation of Gk.\ χολ-ωτός,\iw{AG}{χολωτός} Lat.\ \textit{aegr-ōtus}\iw{Latin}{\textit{aegrōtus}} etc. from ‘deinstrumental’ \iwi{PIE}{*-\textit{oh\textsubscript{1}-to}-} is open to criticism, cf.\ \citet[82--90]{Fortson2020}. For Balto-Slavic, (post-)PIE \iwi{PIE}{*-\textit{ō-to}-} would work as well;
    \item while some Baltic denominatives\is{denominative!verb} are susceptible of a deinstrumental gloss ‘be with X’ (e.g., \iwi{Lith}{\textit{dūmúoti}} ‘smoke; get covered with smoke’), this would be forced for most of them;
    \item an alternative view considers Lith.\ \iwi{Lith}{-\textit{úoti}, -\textit{úoja}} an East Baltic offshoot of the Balto-Slavic denominative\is{denominative!verb} type Lith.\ \iwi{Lith}{-\textit{áuti}, -\textit{áuja}} (\ref{sec:BSlauje}), e.g., \citet[352--353]{Vaillant1966}, \citet{Kortlandt1995}. See Villanueva Svensson (\citeyear{VillanuevaSvensson2014auti}, \citeyear{VillanuevaSvensson2015}) for criticism of different aspects of this approach.
\end{itemize}

\subsection{Bl.-Sl.\ *-\textit{\underline{au̯}-i̯e/a}- (inf.\ *-\textit{\underline{au̯}-tī}, aor.\ *-\textit{au̯-ā}- [!])}
\label{sec:BSlauje}

This suffix is productive\is{productivity} in Lithuanian and Old Prussian: Lith.\ \textit{draug-áuti}, -\textit{áuja}, -\textit{ãvo}\iw{Lith}{\textit{draugáuti, áuja, -ãvo}} ‘be friends’ ($\leftarrow$ \iwi{Lith}{\textit{draũgas}} ‘friend’),\iw{OPr}{\textit{dīnkaut}, 1pl. \textit{dīnkauimai}, pret. \textit{dīnkauts, dinkowats} I} OPr.\ \textit{dīnk-aut}, pres.\ 1pl.\ \textit{dīnk-au(i)mai}, pret.\ \textit{dīnk-auts} (\textit{dink-owats} I) ‘thank’ ($\leftarrow$ Pol. \iwi{Pol}{\textit{dziękować}}). In Lithuanian it derives verbs from nouns of all stems with a basic meaning ‘be X’ (whence ‘work as X’, ‘behave as X’, ‘be engaged in X’, ‘show up as X’, etc.). Deadjectives\is{deadjectival!verb} are rarer, but not exceptional (\iwi{Lith}{\textit{šykštùs}} ‘stingy’ $\rightarrow$ \textit{šýkštauti}, -\textit{auja}\iw{Lith}{\textit{šýkštauti, -auja}} ‘be stingy’). Deverbative\is{deverbal/deverbative!verb} iteratives\is{iterative} also occur, but are generally regarded as secondary\is{derivation!stem-based} (\textit{šaũkti}, -\textit{ia}\iw{Lith}{\textit{šaũkti, -ia}} ‘shout’ $\rightarrow$ \textit{šū́kauti}, -\textit{auja}\iw{Lith}{\textit{šū́kauti, -auja}} ‘whoop’). This type has been lost in Latvian, merging with the verbs in \iwi{Latv}{-\textit{uôt}, -\textit{uõju}} (\ref{sec:BSloje}). The types *-\textit{\underline{au̯}-tī}, *-\textit{\underline{au̯}-i̯e/a}-\iw{PBSl}{*-\textit{au̯-tī}, *-\textit{au̯-i̯e/a}-} and *-\textit{\underline{ō}-tī}, *-\textit{\underline{ō}-i̯e/a}-\iw{PBSl}{*-\textit{ō-tī}, *-\textit{ō-i̯e/a}-} often overlap in Lithuanian as well.

In Slavic, this type is very productive\is{productivity} for both denominatives\is{denominative!verb} and deadjectives:\is{deadjectival!verb}  OCS \textit{dar-ovati}, -\textit{ujǫ}, -\textit{ovaxъ}\iw{OCS}{\textit{darovati, -ujǫ, -ovaxъ}} ‘make a present’ ($\leftarrow$ \iwi{OCS}{\textit{darъ}} ‘present’), \textit{mil-ovati}, -\textit{ujǫ}\iw{OCS}{\textit{milovati, -ujǫ}} ‘have mercy’ ($\leftarrow$ \iwi{OCS}{\textit{milъ}} ‘pitiful; loved’). In addition, it is abundantly used for borrowings (OCS \textit{kup-ovati}, -\textit{ujǫ}\iw{OCS}{\textit{kupovati, -ujǫ}} ‘buy’, cf.\ German \iwi{German}{\textit{kaufen}}). Its use for derived \isi{imperfective}s can be tracked in the recorded history of Old Church Slavonic, cf.\ \citet[29--33]{Schuyt1990}. 

For Balto-Slavic we can confidently reconstruct a denominative suffix\is{denominative!suffix} *-\textit{\underline{au̯}\nobreakdash-tī}, *-\textit{\underline{au̯}\nobreakdash-i̯e/a}-.\iw{PBSl}{*-\textit{au̯-tī}, *-\textit{au̯-i̯e/a}-} It is uncertain whether it was also widely used for deadjectives.\is{deadjectival!verb} From a morphological point of view, it is the only denominative suffix\is{denominative!suffix} with an \textit{ā}-aorist *-\textit{au̯-ā}-,\iw{PBSl}{*-\textit{au̯-ā}- (aor.)} as borne out by the second stem in -\textit{a}- of Slavic -\textit{ov-a-ti}, -\textit{ov-a-xъ}\iw{OCS}{-\textit{ovati, -ovaxъ}}. This is potentially important from a historical perspective, as the Balto-Slavic \textit{ā}-aorist was otherwise limited to primary verbs,\is{derivation!deradical} cf.\ \citet[382--389]{VillanuevaSvensson2020}. To judge by Slavic (the Baltic \textit{ā}-preterit is fully ambiguous) all other denominative suffixes\is{denominative!suffix} had a sigmatic aorist in Balto-Slavic (*-\textit{\underline{ā}-s}-,\iw{PBSl}{*-\textit{ā-s}- (aor.)} *-\textit{\underline{ē}-s}-\iw{PBSl}{*-\textit{ē-s}- (aor.)} etc.).

The origin of Bl.-Sl.\ *-\textit{\underline{au̯}-tī}, *-\textit{\underline{au̯}-i̯e/a}-\iw{PBSl}{*-\textit{au̯-tī}, *-\textit{au̯-i̯e/a}-} is a traditional \textit{locus desperatus} of Balto-Slavic historical grammar, the only certain thing being that it must be an innovation of this branch. I refer to \citet{VillanuevaSvensson2014auti} for a full treatment, including criticism of the widespread idea that it originated in denominatives\is{denominative!verb} from \textit{u}-stem nouns. In this article, I have proposed deriving this suffix from original compounds \is{compounding} with the root \iwi{PIE}{*\textit{h\textsubscript{2}eu̯h\textsubscript{1}}-} ‘help’ (Ved.\ \iwi{Sanskrit}{\textit{ávati}} ‘id.’ etc.). Parallels for such a \isi{reanalysis} include Lat.\ \iwi{Latin}{-\textit{ī̆gāre}}, \iwi{Latin}{-\textit{cināre}} and, especially, OIr.\ -\textit{(a)igithir}. \iw{OIr}{-\textit{aigithir}}

\subsection{Bl.-Sl.\ *-\textit{ī}- (inf.\ *-\textit{\underline{ī}-tī}, aor.\ *-\textit{\underline{ī}-s}- [\textit{vel sim}.])}
\label{sec:BSli}

As a result of some recent insights this suffix is now better understood than just a decade ago. I will therefore treat it in more detail and, more importantly, in a different way than the other denominative formations.\is{denominative!verb}

The origin of Slavic \textit{i}-verbs like OCS \iwi{OCS}{\textit{prositi}, \textit{prosi}-} ‘ask for’ is perfectly clear: they continue PIE \textit{o}-grade \isi{causative}s and iteratives\is{iterative} in \iwi{PIE}{*-\textit{éi̯e/o}-} (PIE \iwi{PIE}{*\textit{mon-éi̯e/o}-} ‘remind’ > Ved.\ \iwi{Sanskrit}{\textit{mānáyati}}, Lat.\ \iwi{Latin}{\textit{moneō}, -\textit{ēre}}) and denominatives\is{denominative!verb} of \textit{o}-stem nominals in \iwi{PIE}{*-\textit{e-i̯é/ó}-} (Ved.\ \iwi{Sanskrit}{\textit{devá}-} ‘god’ $\rightarrow$ \iwi{Sanskrit}{\textit{devayáti}} ‘worship the gods’). The PIE functions are preserved nearly intact in Slavic, where \textit{i}-verbs build iteratives\is{iterative} (OCS \iwi{OCS}{\textit{nesti}, \textit{nesǫ}} ‘carry’ $\rightarrow$ \iwi{OCS}{\textit{nositi}, \textit{nošǫ}} ‘id. (iter.)’), \isi{causative}s (OCS \iwi{OCS}{\textit{tešti}, \textit{tekǫ}} ‘run, flow’ $\rightarrow$ \iwi{OCS}{\textit{točiti}, \textit{točǫ}} ‘let run, let flow’), and denominatives\is{denominative!verb} from both nouns and adjectives: OCS \textit{měr-iti}, -\textit{jǫ}\iw{OCS}{\textit{měriti, -jǫ}} ‘measure’ ($\leftarrow$ \iwi{OCS}{\textit{měra}} ‘measure’), \textit{běl-iti}, -\textit{jǫ}\iw{OCS}{\textit{běliti, -jǫ}} ‘whiten’ ($\leftarrow$ \iwi{OCS}{\textit{bělъ}} ‘white’). \textit{i}-denominatives\is{denominative!verb} are immensely productive.\is{productivity} Deadjectives \is{deadjectival!verb} often (but not exclusively) display \isi{factitive} value. From an accentological point of view, this is the only Slavic denominative suffix\is{denominative!suffix} with an old contrast of tone between present (with non-acute *-\textit{ī}-, CSl.\ 2pl.\ *\textit{nòsite} < *\textit{nosȋte} < Bl.-Sl.\ *\textit{na̍śīte}) and aorist-infinitive (with acute *-\textit{\underline{ī}}-,\iw{PBSl}{*-\textit{ī}- (acute)} CSl.\ inf.\ *\textit{nos\H{i}ti}\iw{PSl}{*\textit{nos\H{i}ti}, 2pl. *\textit{nòsite} < *\textit{nosȋte}} < Bl.-Sl.\ *\textit{na̍ś\underline{ī}tī}\iw{PBSl}{*\textit{naśītī}, 2pl. *\textit{naśīte}}).

From a historical perspective, the Slavic \textit{i}-verbs pose two major problems:

\begin{enumerate}
\item Can Sl.\ \iwi{OCS}{-\textit{i}-} regularly derive from PIE *-\textit{ei̯e/o}-?\iw{PIE}{*-\textit{éi̯e/o}-} If not, how is Sl.\ \iwi{OCS}{-\textit{i}-} to be explained?

\item Is the acute \iwi{PSl}{*-\textit{\H{i}}-} of the aorist-infinitive stem the same as the non-acute \iwi{PSl}{*-\textit{ȋ}-} of the present stem? If not, how is Sl.\ \iwi{PSl}{*-\textit{\H{i}}-} to be explained? Note that Baltic shows aor.-inf.\ \iwi{PSl}{*-\textit{\H{i}}-} to be old with certainty (Lith.\ inf.\ \iwi{Lith}{\textit{prašýti}} = Sl.\ *\textit{pros\H{i}ti}\iw{PSl}{*\textit{pros\H{i}ti}, *\textit{prosȋtь}}).

\end{enumerate}

The second question is rarely posited in this way and, accordingly, has received relatively little attention in the literature.\footnote{See \citet[3]{Klingenschmitt1978}, \citet[207--208]{Jasanoff2017BSl} for scenarios not deriving Bl.-Sl.\ aor.-inf.\ *-\textit{\underline{ī}}-\iw{PBSl}{*-\textit{ī}- (acute)} from PIE *-\textit{ei̯e}-.\iw{PIE}{*-\textit{éi̯e/o}-}} As for the first one, it is probably safe to say that the matter has always been seen with a certain degree of perplexity. Accounts not operating with a regular development PIE *\nobreakdash-\textit{ei̯e}-\iw{PIE}{*-\textit{éi̯e/o}-} > Sl.\ \iwi{OCS}{-\textit{i}-} have usually taken one of the following four options:

\begin{enumerate}
    \item to start from a PIE semithematic or athematic “\textit{ī̆}-present”;\footnote{“Semithematic \textit{i}-presents” were popular in older says of our discipline, see especially \citet[23--29]{Stang1942}. See \citet{Schrijver2003athem} for a modern defense of the PIE “athematic \textit{i}-presents” and \citet[152--155]{Weiss2012} for a different explanation of the Italic and Celtic facts adduced by Schrijver.}
    \item to assume that the -\textit{i}- of \iwi{OCS}{\textit{prositi}, \textit{prosi}-} was taken from that of the \isi{stative} type \iwi{OCS}{\textit{mьněti}, \textit{mьni}-} and/or from \textit{i}-stem denominatives\is{denominative!verb} (PIE “\iwi{PIE}{*-\textit{ī̆-i̯e/o}-}”);\footnote{Both statements are frequent in Slavic handbooks, e.g., \citet[439--440]{Vaillant1966}, \citet[91]{Birnbaum1997}. As often noted (e.g., \citealt[102--103]{Jasanoff1978}), if one type imposed its morphology on the other(s), it is likely that this was the immensely productive\is{productivity} type \textit{prositi}, \textit{prosi}-.}
    \item to postulate complex analogical scenarios;\footnote{E.g., \citet{Kurylowicz1973}, \citet[3--5]{Klingenschmitt1978}.}
    \item to operate with irregular phonological changes.\footnote{E.g., \citet[260: syncope of the second vowel of *-\textit{ei̯e}-]{Matasovic2008}, \citet[60--69: Late Common Slavic contraction across \textit{jod}.]{Andersen2014}.}
\end{enumerate}

Viewed against this background, the assumption of a regular development PIE *-\textit{ei̯e}-\iw{PIE}{*-\textit{éi̯e/o}-} > Sl.\ \iwi{OCS}{-\textit{i}-} becomes, in my view, even more attractive (\emph{all} alternatives are unsatisfactory). The idea is far from new,\footnote{Cf.\ \citet[61--62; 78--79]{Specht1934}, with references to J.\ Schmidt, Bezzenberger, Sommer and Pedersen.} but has certainly gained momentum in recent years; see especially \citet{Hock1995}, \citet[214--222]{Hill2016}, \citet[74--76, 191, 194--199]{VillanuevaSvensson2023a}. A detailed treatment is impossible here, but I would like to stress two insights gained from these works:

\begin{enumerate}
    \item The amount of positive evidence supporting a \textit{regular} development PIE *\nobreakdash-\textit{ei̯e}-\iw{PIE}{*-\textit{éi̯e/o}-} > Bl.-Sl.\ *-\textit{ī}-\iw{PBSl}{*-\textit{ī}- (non-acute)} > Sl.\ \iwi{OCS}{-\textit{i}-} has increased (e.g., Sl.\ inf.\ -\textit{ti},\iw{OCS}{-\textit{ti} (Inf.)} Lith.\ -\textit{ti}\iw{Lith}{-\textit{ti} (Inf.)} < PIE \iwi{PIE}{*-\textit{tei̯ei̯}});

    \item The development PIE *-\textit{ei̯e}-\iw{PIE}{*-\textit{éi̯e/o}-} > Bl.-Sl.\ *-\textit{ī}-\iw{PBSl}{*-\textit{ī}- (non-acute)} was Balto-Slavic in date, as most probably the analogical replacement of *-\textit{ei̯o}- > \iwi{PBSl}{*-\textit{ii̯a}-} with \iw{PBSl}{*-\textit{ī}- (non-acute)}*-\textit{ī}-.\footnote{Against a relatively widespread assumption, the regular reflex of PIE *-\textit{ei̯e/o}-\iw{PIE}{*-\textit{éi̯e/o}-} in Baltic was not (Lith.) \iwi{Lith}{-\textit{ia}} or \iwi{Lith}{-\textit{ija}} (\textit{pace} \citealt[107]{Kortlandt1989}, \citealt[28]{Ostrowski2006}, or \citealt[207]{Jasanoff2017BSl}, among others). See below in the text.} In the aorist-infinitive stem *-\textit{ī}-\iw{PBSl}{*-\textit{ī}- (non-acute)} was secondarily acuted to *-\textit{\underline{ī}}-\iw{PBSl}{*-\textit{ī}- (acute)} on the model of all verbs with a dissyllabic second stem (*-\textit{\underline{ā}}-,\iw{PBSl}{*-\textit{ā}- (acute)} *-\textit{\underline{ē}}-\iw{PBSl}{*-\textit{ē}- (acute)} etc. < *-\textit{ah\textsubscript{2}}-, *-\textit{eh\textsubscript{1}}-).\iw{PIE}{*-\textit{éh\textsubscript{1}-/-h\textsubscript{1}}-}
\end{enumerate}

\tabref{tab:ejeoBSl} gives the development of the PIE *-\textit{ei̯e/o}-presents\iw{PIE}{*-\textit{éi̯e/o}-} in Balto-Slavic.\footnote{Adapted from \citet[198]{VillanuevaSvensson2023a}. Arrows indicate analogical forms. The endings, which are sometimes insecure, cannot be discussed here. Column `Bl. (1)' is deliberately left unaccented as a way to indicate that no direct connection is postulated between PIE stress position (\iwi{PIE}{*\textit{pro\kinvbreve -éi̯e-ti}}) and that of late Balto-Slavic (*\textit{pra̍śīt(i)}).\iw{PBSl}{*\textit{praśīti} (3sg.)}}


\begin{table}
\fittable{
\caption{PIE *-\textit{ei̯e/o}-presents in Balto-Slavic}
\label{tab:ejeoBSl}
\begin{tabular}{lllllll}
\lsptoprule
      & PIE & Bl.-Sl.\ (1) & Bl.-Sl.\ (2) & Sl.\ & Bl.\ (1) & Bl.\ (2) \\
\midrule
 1sg. & *pro\kinvbreve éi̯oh\textsubscript{2} & *praśii̯\underline{ō} & $\rightarrow$ *pra̍śi̯\underline{ō} & *prošǫ̍ & *praši̯\underline{ō} & *praši̯\underline{ō}\\
 2sg. & *pro\kinvbreve éi̯esi & *praśīsi & *pra̍śīš(i) & *prosȋši &  ? &  ?\\
 3sg. & *pro\kinvbreve éi̯eti\iw{PIE}{*\textit{pro\kinvbreve -éi̯e-ti}} & *praśīti\iw{PBSl}{*\textit{praśīti} (3sg.)} & *pra̍śīt(i) & *prosȋtь\iw{PSl}{*\textit{pros\H{i}ti}, *\textit{prosȋtь}} & *praši\iw{PB}{*\textit{praši}} & *praši\\
 1pl. & *pro\kinvbreve éi̯omos & *praśii̯amas & $\rightarrow$ *pra̍śīmas & *prosȋmъ & *prašīme & $\rightarrow$ *prašime\\
 2pl. & *pro\kinvbreve éi̯ete & *praśīte & *pra̍śīte & *prosȋte & *prašīte & $\rightarrow$ *prašite\\
 3pl. & *pro\kinvbreve éi̯onti & *praśii̯anti & $\rightarrow$ *pra̍śīnt(i) & *prosę̑tь &  & \\
{Inf.} &  & *praśītī & $\rightarrow$ *pra̍ś\underline{ī}tī & *pros\H{i}ti & *praš\underline{ī}ti & *praš\underline{ī}ti\\
{Aor.}  & & *praśī-s- & $\rightarrow$ *pra̍ś\underline{ī}-š- & *pros\H{i}xъ & $\rightarrow$ *prašii̯ā & *prašē\\
\lspbottomrule
\end{tabular}}\end{table}


The early Baltic development is the product of a series of important developments:

\begin{itemize}
    \item generalization of the apocopated PIE 3sg.\ *-\textit{ti}\iw{PIE}{*-\textit{ti} (3sg.)} > Bl.-Sl.\ *-\textit{t}\iw{PBSl}{*-\textit{t} (3sg.)} > *-Ø in non-athematic paradigms, yielding *\textit{praš-ī(-t)};\iw{PBSl}{*\textit{praśīti} (3sg.)}
    \item shortening of unstressed non-acute word-final *-\textit{ī}(C),\iw{PBSl}{*-\textit{ī}- (non-acute)} yielding *\textit{praš-i}\iw{PB}{*\textit{praši}} under ‘Bl.\ (1)’ in \tabref{tab:ejeoBSl} (\citealt[214--222]{Hill2016}, \citealt[74--76]{VillanuevaSvensson2023a});
    \item loss of number distinctions in the 3\textsuperscript{rd} person, with concomitant restructuring of verbal paradigms as an endingless 3\textsuperscript{rd} person to which 1\textsuperscript{st} and 2\textsuperscript{nd} singular, plural and dual endings are attached (e.g., Lith.\ pres.\ 3sg.\ \iwi{Lith}{\textit{vẽda}} : 1pl.\ \textit{vẽda-me}, 2pl.\ \textit{vẽda-te}). This allows us to expect the analogical reshuffling given under ‘Bl. (2)’ in \tabref{tab:ejeoBSl}.
\end{itemize}

The most surprising feature of the development of the \textit{ī}-verbs in Baltic, however, is that the expected paradigm does not exist. Instead, we find up to five (!) analogical or mixed types:

\begin{enumerate}
    \item Iterative\is{iterative} and \isi{causative} type Lith.\ \textit{laik-ýti}, \textit{laĩk-o}, \textit{laĩk-\.e},\iw{Lith}{\textit{laikýti, laĩko, laĩk\.e}, 1sg. prs. \textit{laikaũ}} Latv.\ \textit{laîc-ît}, \textit{laîk-u}, \textit{laîc-ĩju};\iw{Latv}{\textit{laîcît, laîku, laîcĩju}} OPr.\ \textit{laik-ūt},\footnote{Old Prussian has secondarily generalized \iwi{PB}{*-\textit{ā}-} to the infinitive stem.} \textit{lāik-u}\iw{OPr}{\textit{laikūt, lāiku}} ‘hold’, with \textit{ā}-present (\ref{sec:BSlaje}).

    \item Type OLith.\ \textit{tar-ýti}, \textit{tãr-ia}, \textit{tãr-\.e}\iw{OLith}{\textit{tarýti, tãria, tãr\.e}} ‘say’, attested in relics in Old and dialectal Lithuanian (\citealt[327--328]{Stang1966}). 

    \item Iteratives\is{iterative} with ‘extra -\textit{i}-’ like Lith.\ \textit{lándž-ioti}, -\textit{ioja}, -\textit{iojo}\iw{Lith}{\textit{lándžioti, -ioja, -iojo}} ‘creep, crawl about’, Latv.\ \iwi{Latv}{\textit{luõžât}, -\textit{ãju}} ‘id.’, analyzable as \iwi{PB}{*-\textit{i}-} (pres.\ *\textit{praš-i}\iw{PB}{*\textit{praši}}) + *-\textit{\underline{ā}ti}, *-\textit{\underline{ā}i̯a}.\iw{PBSl}{*-\textit{ā-tī}, *-\textit{ā-i̯e/a}-}

    \item Causative\is{causative} and \isi{factitive} type Lith.\ \textit{dẽg-inti}, -\textit{ina}, -\textit{ino},\iw{Lith}{\textit{dẽginti, -ina, -ino}} Latv.\ \textit{dedz-inât}, -\textit{ina}, \nobreakdash-\textit{inãja}\iw{Latv}{\textit{dedzinât, -ina, -inãja}} ‘burn (tr.\is{transitive})’, OPr.\ \textit{mukint}, pres.\ \textit{mukinna}\iw{OPr}{\textit{mukint, mukinna}} ‘teach’, analyzable as \iwi{PB}{*-\textit{i}-} (pres.\ *\textit{praš-i})\iw{PB}{*\textit{praši}} + \iwi{PBSl}{*-\textit{na}-} (see below \ref{sec:BSlneu}).\footnote{Cf.\ \citet{VillanuevaSvenssonforth}. See \citet[175--176]{Gorbachov2007} for an alternative account that, however, also derives Bl.\ \iwi{PB}{*-\textit{in}-} from the Balto-Slavic \textit{ī}-verbs. The idea is not new, cf.\ \citet[411]{Otrebski1965}, \citet[227--228]{Vaillant1966}, \citet[173]{Schmalstieg2000}. See \citet{Kortlandt1989}, \citet{Hajnal1996} for different views.}

    \item Denominative\is{denominative!verb} type Lith.\ \textit{dal-ýti}, -\textit{ìja}, -\textit{ìjo}\iw{Lith}{\textit{dalýti, -ìja, -ìjo}} ‘divide’, Latv.\ \textit{medît}, -\textit{ĩju}\iw{Latv}{\textit{medît, -ĩju}} ‘hunt’, OPr.\ \iwi{OPr}{\textit{madlīt}, \textit{madli}, 1pl.\ \textit{madlimai}} ‘pray’.\footnote{Cf.\ \citet{VillanuevaSvensson2023b}, with arguments against the widespread derivation from \textit{i}-stem denominatives \is{denominative!verb} (e.g., \citealt[367]{Stang1966}, \citealt[128]{Ostrowski2006}, \citealt[279--280]{Pakerys2011}, among others).}

\end{enumerate}

To disentangle the extraordinarily complex development of the Balto-Slavic \textit{ī}-verbs in Baltic is a major task for the future. It is no less important to stress that the Balto-Slavic \textit{ī}-verbs are better understood today than just a decade ago, to a large degree thanks to recent advances in historical phonology. As for their position in the Balto-Slavic denominative system,\is{denominative!verb} probably the main question is whether they were productively\is{productivity} used to build deadjectival \is{deadjectival!verb} \isi{factitive}s in addition to plain (\isi{transitive}) denominatives; see below (\ref{sec:BSlneu}).

\subsection{Bl.-Sl.\ *-\textit{\underline{ē}-i̯e/a}- (inf.\ *-\textit{\underline{ē}-tī}, aor.\ *-\textit{\underline{ē}-s}- [\textit{vel sim}.])}
\label{sec:BSleje}

This suffix is productive\is{productivity} in all three Baltic languages, e.g., Lith.\ \textit{stor-\dotacute{e}ti}, \mbox{-\textit{\dotacute{e}ja},} \mbox{-\textit{\dotacute{e}jo}}\iw{Lith}{\textit{stor\dotacute{e}ti, -\dotacute{e}ja, -\dotacute{e}jo}} ‘grow fat’ ($\leftarrow$ \iwi{Lith}{\textit{stóras}} ‘fat’); Latv.\ \textit{me\ltilded n-êt}, -\textit{ẽju}, -\textit{ẽju}\iw{Latv}{\textit{me\ltilded nêt, -ẽju, -ẽju}} ‘become black’ ($\leftarrow$ \mbox{\iwi{Latv}{\textit{mȩ\ltilded ns}}} ‘black’); OPr.\ \textit{druw-īt}, \textit{druw-ē},  1pl.\ \textit{druw-ēmai}\iw{OPr}{\textit{druwīt, druwē}, 1pl. \textit{druwēmai}} ‘believe’ ($\leftarrow$ \iwi{OPr}{\textit{druwis}} ‘faith’). It builds \isi{fientive}s from adjectives \is{deadjectival!verb} and, less often, nouns. Essives\is{essive} occur, but only as relics: Lith.\ \iwi{Lith}{\textit{gars\dotacute{e}ti}, -\textit{\dotacute{e}ja}} ‘be famed’ ($\leftarrow$ \iwi{Lith}{\textit{garsùs}} ‘famous’), Latv.\ \iwi{Latv}{\textit{klusêt}, -\textit{ẽju}} ‘be silent’ ($\leftarrow$ \iwi{Latv}{\textit{kluss}} ‘silent’); cf.\ \citet{Pakerys2006}, \citet[109--115]{Ostrowski2006}. The decay of the \isi{essive} value is probably related to the decay of nasal \isi{fientive}s \is{nasal present}  (\ref{sec:BSlgrun}) and the expansion of \isi{essive}s in -\textit{áuti}\iw{Lith}{-\textit{áuti}, -\textit{áuja}} (\ref{sec:BSlauje}) and -\textit{úoti}\iw{Lith}{-\textit{úoti}, -\textit{úoja}} (\ref{sec:BSloje}). Fientive\is{fientive} value may have begun with preverbs,\is{preverb} as in Slavic.\footnote{Similarly \citet[115]{Ostrowski2006}.} From a morphological point of view, it is well known that Baltic is one of the few branches distinguishing between \textit{ē}-denominatives\is{denominative!verb} (with \mbox{-\textit{\.eja}}-presents) and \textit{ē}-deverbatives\is{deverbal/deverbative!verb} (with presents in -\textit{ti},\iw{Lith}{-\textit{ti} (Prs.)} \iwi{Lith}{-\textit{i}}, \iwi{Lith}{-\textit{a}} and, rarely, \mbox{\iwi{Lith}{-\textit{ia}}}). There is, however, a certain degree of contamination between the different present types. In particular, original denominatives\is{denominative!verb} with \isi{stative} value tend to acquire deverbative\is{deverbal/deverbative!verb} morphology (e.g., Lith.\ \iwi{Lith}{\textit{myl\dotacute{e}ti}, \textit{mýli}}, Latv.\ \iwi{Latv}{\textit{mĩlêt}, -\textit{u}} ‘love’, cf.\ Lith.\ dial.\ \iwi{Lith}{\textit{mýlas}}, Sl.\ \iwi{PSl}{*\textit{m\H{i}lъ}} ‘loved’).

Slavic essentially agrees with Baltic. The suffix \iwi{PSl}{*-\textit{\HighH{ě}ti}, -\textit{\HighH{ě}jǫ}} is productive\is{productivity} and derives \isi{fientive}s and, more rarely, \isi{essive}s from adjectives \is{deadjectival!verb}  and, less often,\is{denominal!verb} nouns: \textit{cěl-ěti}, -\textit{ějǫ}\iw{OCS}{\textit{cělěti, -ějǫ}} ‘recover’ ($\leftarrow$ \iwi{OCS}{\textit{cělъ}} ‘healthy’), \textit{um-ěti}, -\textit{ějǫ}\iw{OCS}{\textit{uměti, -ějǫ}} ‘be able’ ($\leftarrow$ \iwi{OCS}{\textit{umъ}} ‘mind’).\footnote{\citet[366--373]{Vaillant1966}.} According to \citet[369]{Vaillant1966} the \isi{fientive} value started in preverbed\is{preverb} \isi{perfective}s: OCS \iwi{OCS}{\textit{bogatěti}} ‘be rich’ : \textit{raz-bogatěti}\iw{OCS}{\textit{razbogatěti}} ‘become rich’ ($\leftarrow$ \iwi{OCS}{\textit{bogatъ}} ‘rich’; cf.\ \isi{factitive} \iwi{OCS}{\textit{bogatiti}} ‘make rich’). It clearly predominates already in Old Church Slavonic.

For Balto-Slavic, we can safely reconstruct a deadjective \is{deadjectival!verb} suffix *-\textit{\underline{ē}-tī}, *-\textit{\underline{ē}-i̯e/a}-.\iw{PBSl}{*-\textit{ē-tī}, *-\textit{ē-i̯e/a}-} It is hard to say whether it also built desubstantives. I assume that it regularly built \isi{essive}s, even though this is mostly based on internal reconstruction. Fientive\is{fientive} value is far better attested in historical Baltic and Slavic, but is easily explained as parallel innovations.

An important topic that can only be tangentially mentioned here is the role of the second stem *-\textit{\underline{ē}}-\iw{PBSl}{*-\textit{ē}- (acute)} in deverbatives.\is{deverbal/deverbative!verb} This is of course not exclusive of Balto-Slavic, but this branch, especially Baltic, stands out by the richness of deverbative\is{deverbal/deverbative!verb} formations:

\begin{enumerate}
    \item With \isi{stative} \textit{i}-presents: Lith.\ \iwi{Lith}{\textit{bud\dotacute{e}ti}, \textit{bùdi}}, OCS \textit{bъděti}, \textit{bъdi}-\iw{OCS}{\textit{bŭděti}, \textit{bŭdi}-} ‘be awake’. This is the best-known type.\footnote{See \citet{Jasanoff2004} and, more recently, \citet{Yakubovich2014}, \citet{Grestenberger2019} on the origin of the Balto-Slavic \textit{i}-presents. See \citet[109--110]{Kortlandt1989}, \citet{Hardarson1998}, \citet[216--219]{Hill2016} for alternative views.}
    \item With athematic\is{root!present} present: OLith.\ \iwi{OLith}{\textit{gélb\.eti}, -\textit{mi}} ‘save’, OCS \iwi{OCS}{\textit{věděti}, \textit{vědě}/\textit{věmь}} ‘know’. This type is only well represented in Old Lithuanian.
    \item With thematic present\is{thematic!present}: Lith.\ \iwi{Lith}{\textit{tek\dotacute{e}ti}, \textit{tẽka}} ‘run, flow’. Only in Baltic, but with potential relics in Slavic.\footnote{Cf.\ \citet{Koelln1977}, \citet[470--473]{VillanuevaSvensson2017}. The Lithuanian evidence is collected in \citet{Jakulis2004}.}
    \item With \textit{ā}-present: OCS \iwi{OCS}{\textit{iměti}, \textit{imamь}} ‘have’ (only certain example).\footnote{Old Prussian verbs like \iwi{OPr}{\textit{billītwei}, \textit{billē} {\textasciitilde} \textit{billā}} ‘speak’, with chaotic variation between present \iwi{OPr}{-\textit{ā}} and \iwi{OPr}{-\textit{ē}}, could belong here as well (cf.\ \citealt[187--204]{Kaukiene2004}).}
    \item Several derived verbal classes in East Baltic:
    \begin{enumerate}
        \item \isi{durative}s with acute metatony (\iwi{Lith}{\textit{stéip\.eti}, -\textit{\.eja}} ‘lie dying’ $\leftarrow$ \iwi{Lith}{\textit{stìpti}} ‘die (of animals)’);
        \item ‘\isi{intensive}s’, also with acute metatony (\iwi{Lith}{\textit{svérd\.eti}, -\textit{i}} ‘stagger’ $\leftarrow$ \iwi{Lith}{\textit{svìrti}} ‘bend down’); 
        \item diminutives with non-acute lengthened grade (\textit{pa-\.ej\dotacute{e}ti}, -\textit{ė̃ji}/-\textit{\.ej\dotacute{e}ja}\iw{Lith}{\textit{pa\.ej\dotacute{e}ti, -ė̃ji/-\.ej\dotacute{e}ja}} ‘go a little’ $\leftarrow$ \iwi{Lith}{\textit{eĩti}} ‘go’); 
        \item \isi{causative}s (only in Latvian, \iwi{Latv}{\textit{aûdzêt}, -\textit{ẽju}} ‘cultivate, grow (tr.\is{transitive})’ $\leftarrow$ \iwi{Latv}{\textit{aûgt}} ‘grow (intr.\is{intransitive})’); 
        \item several subtypes with ‘extended’ suffix like \iwi{Lith}{-\textit{al\.eti}}, \iwi{Lith}{-\textit{el\.eti}}, \iwi{Lith}{-\textit{er\.eti}}; \iwi{Lith}{-\textit{tel(\.e)ti}}, \iwi{Lith}{-\textit{ter(\.e)ti}}; \iwi{Lith}{-\textit{in\.eti}}; \iwi{Lith}{-\textit{s\.eti}}.\footnote{See \citet[11--32]{Ostrowski2006}, \citet{Serzant2014}, \citet[100--114]{RothsteinDowden2022} for different issues related to this material.}
        \end{enumerate}
\end{enumerate}

The precise prehistory of some of these formations is still unclear, just as their interplay with the \textit{ē}-denominatives.\is{denominative!verb} As for the latter, Balto-Slavic deadjectives\is{deadjectival!verb} in *-\textit{\underline{ē}i̯e/a}-\iw{PBSl}{*-\textit{ē-tī}, *-\textit{ē-i̯e/a}-} are generally regarded as an archaism (PIE \iwi{PIE}{*-\textit{eh\textsubscript{1}-i̯e/o}-}). It is important to note that reflexes of *-\textit{eh\textsubscript{1}(-i̯e/o)}-\iw{PIE}{*-\textit{eh\textsubscript{1}-i̯e/o}-} are exclusively denominative\is{denominative!suffix} in Anatolian and Indo-Iranian (cf.\ \citealt{Yakubovich2014}). This strongly suggests that the presence of *-\textit{eh\textsubscript{1}(-i̯e/o)}-\iw{PIE}{*-\textit{eh\textsubscript{1}-i̯e/o}-} among deverbatives\is{deverbal/deverbative!verb} is post-Indo-European. With all due caution, I suggest viewing the traditional deverbative\is{deverbal/deverbative!verb} \textit{ē}-statives\is{stative} as an innovation of the North-Western language area (the path taken by the Greek η-aorist was completely different). The idea, needless to say, requires more work.

\subsection{Bl.-Sl.\ type *\textit{gru-n-b-e/a}- (inf.\ *\textit{grub-tī}, aor.\ *\textit{grub-e/a}-)}
\label{sec:BSlgrun}

As is well known, \isi{nasal present}s with \isi{anticausative} and \isi{inchoative} value constitute one of the defining features of Balto-Slavic and Germanic, e.g., Lith.\ \iwi{Lith}{\textit{lìpti}, \textit{lim̃pa}, \textit{lìpo}} ‘stick to’, OCS \textit{pri-lь(p)nǫti}, -\textit{lь(p)nǫ}, -\textit{lьpъ}\iw{OCS}{\textit{prilĭ(p)nǫti, -lĭ(p)nǫ, -lĭpŭ}} ‘id.’, Go. \textit{af-lifnan}, -\textit{lifniþ}, \nobreakdash-\textit{lifnoda}\iw{Goth}{\textit{aflifnan, -lifniþ, -lifnoda}} ‘be left over’. See \citet{Gorbachov2007} for the reconstruction of the \isi{Northern Indo-European} prototype *\textit{bʰu-n-dʰ-é-ti} ‘awakes’ (aor.\ *\textit{bʰudʰ-é-t}) and \citet{VillanuevaSvensson2011} for my views on their origin. The same formation is productive\is{productivity} for deadjectival \is{deadjectival!verb} \isi{fientive}s in all three branches, more rarely for desubstantives: Go. \iwi{Goth}{\textit{fulls}} ‘full’ $\rightarrow$ \textit{(ga-)fullnan}\iw{Goth}{\textit{(ga)fullnan}} ‘become filled’, Lith.\ \iwi{Lith}{\textit{šlùbas}} ‘lame’ $\rightarrow$ \iwi{Lith}{\textit{šlùbti}, \textit{šlum̃ba}} ‘become lame’, OCS \textit{slěpъ}\iw{OCS}{\textit{slěpŭ}} ‘blind’ $\rightarrow$ \textit{o-slьpnǫti}\iw{OCS}{\textit{oslĭpnǫti, -slĭpnǫ, -slĭpŭ}} ‘go blind’.

Fientives\is{fientive} are very well represented in East Baltic: Lith.\ \iwi{Lith}{\textit{šlàpti}, \textit{šlam̃pa}, \textit{šlãpo}} ‘become wet’ ($\leftarrow$ \iwi{Lith}{\textit{šlãpias}} ‘wet’), Latv.\ \iwi{Latv}{\textit{slapt}, \textit{slùopu}, \textit{slapu}} ‘id.’ ($\leftarrow$ \iwi{Latv}{\textit{slapjš}}).\footnote{The evidence is collected in \citet{Pakalniskiene2000}.} The type, however, is not productive\is{productivity} anymore, its place being taken by \iwi{Lith}{-\textit{\.eti}, -\textit{\.eja}} (\ref{sec:BSleje}). The absence of denominatives\is{denominative!verb} in Old Prussian is surely due to chance.
 
Slavic essentially agrees with Baltic: OCS \textit{slěpъ}\iw{OCS}{\textit{slěpŭ}} ‘blind’ $\rightarrow$ \textit{o-slьpnǫti}\iw{OCS}{\textit{oslĭpnǫti, -slĭpnǫ, -slĭpŭ}} ‘go blind’. Denominatives \is{denominative!verb} are less well represented than in Baltic, no doubt because the \isi{inchoative} type \iwi{OCS}{\textit{bъ(d)nǫti}} itself ceased to be productive\is{productivity} and because verbs in \nobreakdash-\textit{ěti}, \nobreakdash-\textit{ějǫ}\iw{OCS}{-\textit{ěti, -ějǫ}} became productive\is{productivity} for fientives\is{fientive}. An obvious archaism is that zero grade is dominant in the older period: OCS \iwi{OCS}{\textit{gluxъ}} ‘deaf’ $\rightarrow$ \textit{o-glъxnǫti}\iw{OCS}{\textit{oglъxnǫti}} ‘become deaf’, \iwi{OCS}{\textit{xromъ}} ‘lame’ $\rightarrow$  \textit{o-xrъmnǫti}\iw{OCS}{\textit{oxrъmnǫti}} ‘grow lame’ (cf.\ \citealt[241--242]{Vaillant1966}). This feature recurs in Germanic (Gmc. \iwi{PGmc}{*\textit{starka}-} ‘strong’ $\rightarrow$ Go. \textit{ga-staurknan}\iw{Goth}{\textit{gastaurknan}} ‘become rigid’, ON \iwi{ON}{\textit{storkna}} ‘coagulate’) and can be safely projected back into \isi{Northern Indo-European}. Baltic has some relics as well: Lith.\ \iwi{Lith}{\textit{bjùrti}, \textit{bjū̃ra}/\textit{bjùrsta}} ‘become ugly’ ($\leftarrow$ \iwi{Lith}{\textit{bjaurùs}} ‘ugly’).\footnote{Cf.\ \citet{VillanuevaSvensson2016}.}

For Balto-Slavic, we can confidently reconstruct a class of deadjectival\is{deadjectival!verb} \isi{fientive}s with zero grade of the root. The type was identical with the \isi{inchoative} \textit{n}-presents \is{nasal present} and must have involved, among other features, a thematic aorist.\is{thematic!aorist} From a historical point of view, it is clear that the \isi{Northern Indo-European} \isi{fientive}s are an offshoot of the deverbative\is{deverbal/deverbative!verb} \isi{anticausative}-\isi{inchoative}s, obviously via secondary association of some deverbatives\is{deverbal/deverbative!verb} to adjectives from the same root.

\subsection{Bl.-Sl.\ \textit{*plit-ne-ti} (inf.\ \textit{*plit-neu̯-tī}, aor.\ \textit{*plit-neu̯-s}- [\textit{vel sim.}])}
\label{sec:BSlneu}

The last denominative\is{denominative!verb} formation is a recent addition that is likely to remain controversial, at least for some years. I refer to \citet{VillanuevaSvenssonforth} for full argumentation of the ideas presented here.

Balto-Slavic clearly had an elaborated system of deadjectives \is{deadjectival!verb} including morphologically differentiated \isi{essive}s (\ref{sec:BSleje}) and \isi{fientive}s (\ref{sec:BSlgrun}). As for the \isi{factitive}s, probably most scholars would project into Balto-Slavic the Slavic \textit{i}-\isi{factitive} type OCS \iwi{OCS}{\textit{cělъ}} ‘healthy’ $\rightarrow$ \iwi{OCS}{\textit{cěliti}} ‘heal’ (\ref{sec:BSli}). This is clearly an extension of \isi{transitive} desubstantives, no doubt supported by the \isi{causative} value of PIE *-\textit{ei̯e/o}-.\iw{PIE}{*-\textit{éi̯e/o}-} The process may go back to Balto-Slavic, but may equally well be exclusively Slavic.\footnote{The same development took place in Germanic (Go. \iwi{Goth}{\textit{hails}} ‘healthy’ $\rightarrow$ \iwi{Goth}{\textit{hailjan}} ‘heal’) and Sanskrit (\iwi{Sanskrit}{\textit{taruṇa}-} ‘young’ $\rightarrow$ \iwi{Sanskrit}{\textit{taruṇaya}-} ‘make young’).} From a functional point of view, the Slavic \textit{i}-verbs have a clear pendant in the Baltic verbs in *-\textit{\underline{in}ti}, *-\textit{ina},\iw{PB}{*-\textit{inti}, *-\textit{ina}} which are productive\is{productivity} for \isi{factitive}s (Lith.\ \textit{saldùs}\iw{Lith}{\textit{saldùs} m., \textit{saldì} f.} ‘sweet’ $\rightarrow$ \iwi{Lith}{\textit{sáldinti}} ‘sweeten’) and \isi{causative}s (Lith.\ \iwi{Lith}{\textit{mir̃ti}} ‘die’ $\rightarrow$ \iwi{Lith}{\textit{marìnti}} ‘kill’). As noted above (\ref{sec:BSli}), the \iwi{PB}{*-\textit{i}-} of the present *-\textit{ina}\iw{PB}{*-\textit{inti}, *-\textit{ina}} can now be identified with the 3sg.\ *-\textit{i}\iw{PB}{*-\textit{i} (3sg.)} of the Balto-Slavic *-\textit{ī}-verbs in Baltic. The obvious analysis *-\textit{i}+\textit{na}\iw{PB}{*-\textit{inti}, *-\textit{ina}} that this leads to leaves us with a \iwi{PB}{*-\textit{na}} that has no obvious source within Baltic (the subsequent \isi{reanalysis} of *-\textit{i}-\textit{na} as *-\textit{in-a}\iw{PB}{*-\textit{inti}, *-\textit{ina}} is essentially unproblematic).

Let us now turn to PIE. In the 1970s, \citet{Koch1973}, Nussbaum (\textit{apud} \citealt[26]{Tucker1981}) and \citet{Tucker1981} argued that PIE possessed a class of \textit{nu}-\isi{factitive}s derived from \textit{u}-stem adjectives. The type is faithfully continued in the Hittite \isi{factitive}s and \isi{causative}s in -\textit{nu-\textsuperscript{mi}}\iw{Hitt}{-\textit{nu}-} (Hitt.\ \iwi{Hitt}{\textit{daššu}-} ‘strong’ $\rightarrow$ \textit{dašš(a)nu-\textsuperscript{mi}}\iw{Hitt}{\textit{daš(ša)nu}-} ‘strengthen’), in the Greek denominatives\is{denominative!verb}  in \iwi{AG}{-ῡ́νω} (\iwi{AG}{βαθύς} ‘deep’ $\rightarrow$ \iwi{AG}{βαθῡ́νω} ‘deepen’) and, probably, \iwi{AG}{-αίνω} (\iwi{AG}{κῡδρός} ‘glorious’ $\rightarrow$ \iwi{AG}{κῡδαίνω} ‘glorify’),\footnote{The connection of Gk.\ \iwi{AG}{-ῡ́νω} and Hitt.\ {-\textit{nu-\textsuperscript{mi}}}\iw{Hitt}{-\textit{nu}-} was an important aspect of \citegen{Koch1973} argument and has been upheld several times. See \citet{VillanuevaSvensson2024}, building on \citet[22--29]{Tucker1981}, for the idea that Gk.\ \iwi{AG}{-αίνω} continues PIE \textit{nu}-\isi{factitive}s as well.} and, finally, in scattered relics elsewhere in the family (e.g., Ved.\ \iwi{Sanskrit}{\textit{dabhnóti}} ‘deceive’, cf.\ Hitt. \textit{tepnu-\textsuperscript{mi}}\iw{Hitt}{\textit{tepnu}-}  ‘diminish; despise’ $\leftarrow$ \iwi{Hitt}{\textit{tēpu}-} ‘small’). The idea has clearly gained momentum in recent years.\footnote{Cf.\ \citet[206--210]{Steer2013NIP}, \citet{Oettinger2018}, Sasseville (\citeyear[484--486]{Sasseville2021} and this volume), Villanueva Svensson (\citeyear{VillanuevaSvensson2024, VillanuevaSvenssonforth}).} In addition, recent research has refined some of the initial assumptions:

\begin{itemize}
    \item PIE \textit{nu}-\isi{factitive}s were part of the Caland system\is{Caland!system (CS)}  (\citealt[112--120]{Rau2009}, \citealt{Oettinger2018});
    \item new evidence from Luvian suggests that the \textit{nu}-\isi{factitive}s followed the \textit{h\textsubscript{2}}\textit{e}-conjugation, whereas deverbatives\is{deverbal/deverbative!verb} followed the \isi{\textit{mi}-conjugation} (\citealt[485--486]{Sasseville2021}). I explicitly highlight that Greek \isi{factitive}s in \iwi{AG}{-ῡ́νω} and \iwi{AG}{-αίνω} continue PIE \textit{nu}-\isi{factitive}s, as this implies that this type was preserved as a distinct class in Core Indo-European.
\end{itemize}

In Balto-Slavic (as in other branches), a \textit{h\textsubscript{2}}\textit{e}-conjugation suffix is expected to be thematized\is{thematization} at an early date, yielding \iwi{PBSl}{*-\textit{nu̯e/o}-} or \iwi{PBSl}{*-\textit{neu̯e/o}-}. \textit{Qua} secondary verb formation\is{derivation!stem-based} \textit{nu}-\isi{factitive}s lacked an aorist. In Balto-Slavic, this slot is expected to be filled with the present stem minus *-\textit{(i̯)e/o}-, yielding aor.-inf.\ \iwi{PBSl}{*-\textit{nu}-} or \iwi{PBSl}{*-\textit{neu̯}-}. The major claim of \citet{VillanuevaSvenssonforth} is that the Slavic type pres.\ \textit{ri-ne}-, inf.\ \textit{ri-nǫ-ti}, aor.\ \textit{ri-nǫ-xъ},\iw{OCS}{\textit{rinǫti, rine-, rinǫxъ}, ppp \textit{rinovenъ}} past pass. ptcp. \is{participle} \textit{ri-nov-enъ} ‘cast, push’ goes back entirely to the Balto-Slavic descendant of the PIE \textit{nu}-\isi{factitive}s.\is{nasal present} The aorist-infinitive stem extended full-grade \iwi{PBSl}{*-\textit{neu̯}-}. In the present, zero-grade \iwi{PBSl}{*-\textit{nu̯e/o}-} was simplified to *\nobreakdash-\textit{ne/o}-\iw{PBSl}{*-\textit{ne/o}-} by regular sound change. \textit{nu}-\isi{factitive}s \is{nasal present} were lost in the prehistory of Slavic, being replaced by the \textit{ī}-verbs (\ref{sec:BSli}). When the Slavic aspect system arose, remaining \textit{nu}-\isi{factitive}s \is{nasal present} were pressed into new service as secondary \isi{perfective}s.\is{derivation!stem-based} In Baltic, on the other hand, \textit{nu}-\isi{factitive}s \is{nasal present} were contaminated with \textit{ī}-verbs (*\nobreakdash-\textit{i}+\textit{na}), finally leading to the productive\is{productivity} \isi{factitive} and \isi{causative} type *-\textit{\underline{in}ti}, *-\textit{ina}.\iw{PB}{*-\textit{inti}, *-\textit{ina}}

If this is correct, it implies that Balto-Slavic had a distinct formation for deadjective \is{deadjectival!verb} \isi{factitive}s that directly continues the PIE \textit{nu}-\isi{factitive}s.\is{nasal present}

\section{Balto-Slavic denominatives: some general issues}
\label{sec:BSlissues}
\subsection{Introduction}
\label{sec:BSlissuesintro}

The dossier of Balto-Slavic denominatives\is{denominative!verb}  included the suffixes *-\textit{\underline{ā}-i̯e/a}-,\iw{PBSl}{*-\textit{ā-tī}, *-\textit{ā-i̯e/a}-} \mbox{*-\textit{\underline{ō}-i̯e/a}-},\iw{PBSl}{*-\textit{ō-tī}, *-\textit{ō-i̯e/a}-} \mbox{*-\textit{\underline{au̯}-i̯e/a}-},\iw{PBSl}{*-\textit{au̯-tī}, *-\textit{au̯-i̯e/a}-} *-\textit{ī}-\iw{PBSl}{*-\textit{ī}- (non-acute)} (< PIE *-\textit{ei̯e/o}-)\iw{PIE}{*-\textit{éi̯e/o}-} and a clearly delimited deadjectival \is{deadjectival!verb} system involving \isi{essive}s (*\textit{grau̯b-\underline{ē}-tī}, *\textit{grau̯b-\underline{ē}-i̯e/a}-\iw{PBSl}{*\textit{grau̯bētī}, *\textit{grau̯bēi̯e/a}-} ‘be coarse’), \isi{fientive}s (*\textit{grub-tī}, \mbox{*\textit{gru-m-be/a}-}\iw{PBSl}{*\textit{grubtī}, *\textit{grumbe/a}-} ‘become coarse’) and, probably, \isi{factitive}s (*\textit{grau̯b-n\underline{au̯}-tī}, *\textit{grau̯b-ne/a}-\iw{PBSl}{*\textit{grau̯bnau̯tī}, *\textit{grau̯bne/a}-} ‘make coarse’). The \isi{factitive} is the only (potentially) controversial slot (other authors would reconstruct *\textit{grau̯b-\underline{ī}-tī}, *\textit{grau̯b-ī}-).\iw{PBSl}{*\textit{grau̯bītī}, *\textit{grau̯bī}-} Needless to say, the boundaries between deadjectives \is{deadjectival!verb} and denominatives\is{denominative!verb}  are unlikely to have been very strict.

This survey does not answer all questions related to the reconstruction of the Balto-Slavic denominative\is{denominative!verb} system. Thus, the distribution of the different denominative suffixes\is{denominative!suffix} is not fully clear (apart from the connection of *-\textit{\underline{ā}-i̯e/a}-\iw{PBSl}{*-\textit{ā-tī}, *-\textit{ā-i̯e/a}-} and *\nobreakdash-\textit{\underline{ō}\nobreakdash-i̯e/a}-\iw{PBSl}{*-\textit{ō-tī}, *-\textit{ō-i̯e/a}-} to \textit{ā}- and \textit{o}-stems, respectively), and the same holds true for their semantics. Problems of this sort are part and parcel of the reconstruction of denominatives,\is{denominative!verb} but further research could bring in some clarity (one must note, at this point, that study of Baltic and Slavic denominatives\is{denominative!verb}  is almost never carried out within a proper Balto-Slavic perspective). Other issues constitute global problems of the Balto-Slavic verb. Thus, although the formation of the Balto-Slavic second stem is generally clear, the way it came into being is not well understood. The caveat with which I accompany aorists like *-\textit{\underline{ā}-s}-\iw{PBSl}{*-\textit{ā-s}- (aor.)} (‘\textit{vel sim}.’), for instance, indicates that our reconstruction of the Balto-Slavic aorist relies almost exclusively on Slavic and is, accordingly, less certain than other forms. Again, future research may bring in some clarity on this and related issues. In what follows I address some general issues of particular interest for Indo-European studies.

\subsection{PIE background 1: The contribution of Balto-Slavic}
\label{sec:BSlPIEcontribution}

From a historical perspective, the Balto-Slavic denominative suffix\is{denominative!suffix} inventory is, as expected, a mix of preservation, loss, and innovation. Clearly innovated are the \isi{fientive} \isi{nasal present}s and the denominative suffixes\is{denominative!suffix} *-\textit{\underline{au̯}-i̯e/a}-\iw{PBSl}{*-\textit{au̯-tī}, *-\textit{au̯-i̯e/a}-} and, probably, *\nobreakdash-\textit{\underline{ō}-i̯e/a}-.\iw{PBSl}{*-\textit{ō-tī}, *-\textit{ō-i̯e/a}-} These innovations took place at different stages and are not without interest from a broader perspective. If my account of *-\textit{\underline{au̯}-i̯e/a}-\iw{PBSl}{*-\textit{au̯-tī}, *-\textit{au̯-i̯e/a}-} is correct (\ref{sec:BSlauje}), it shows that the creation of denominative suffixes\is{denominative!suffix} out of second compound members\is{compounding} is more widespread than one would otherwise expect. From a PIE perspective *-\textit{\underline{ō}-i̯e/a}-\iw{PBSl}{*-\textit{ō-tī}, *-\textit{ō-i̯e/a}-} is the most interesting item. Following the \textit{communis opinio}, in Section \ref{sec:BSloje} I have assumed that it is a Balto-Slavic innovation. The issue, however, certainly deserves more work. In a more pessimistic vein, the contrast with Gk.\ \iwi{AG}{-όω} illustrates a recurrent problem with the Balto-Slavic evidence. Careful attention to the inflectional, derivational and semantic profile of the -όω-denominatives\is{denominative!suffix} in oldest Greek allowed \citet[15--22]{Tucker1981} to clearly distinguish two different subtypes with, perhaps, two different prehistories. This type of philological work is rarely possible for Balto-Slavic.

As for preservation, the suffixes *-\textit{\underline{ā}-i̯e/a}-\iw{PBSl}{*-\textit{ā-tī}, *-\textit{ā-i̯e/a}-}, *-\textit{\underline{ē}-i̯e/a}-,\iw{PBSl}{*-\textit{ē-tī}, *-\textit{ē-i̯e/a}-} and *-\textit{ī}-\iw{PBSl}{*-\textit{ī}- (non-acute)} (PIE \iwi{PIE}{*-\textit{ah\textsubscript{2}-i̯e/o}-}, *\nobreakdash-\textit{eh\textsubscript{1}-i̯e/o}-,\iw{PIE}{*-\textit{eh\textsubscript{1}-i̯e/o}-} *-\textit{ei̯e/o}-)\iw{PIE}{\textit{*-éi̯e/o}-} are well-known and little needs to be added here. Far more interesting are the \textit{nu}-presents \is{nasal present} (\ref{sec:BSlneu}) and the \textit{ā}-presents (\ref{sec:BSlaje}), if correctly interpreted. The matter can be briefly formulated as follows. The Vedic denominative system\is{denominative!verb}  is extraordinarily simple (nominal stem + \iwi{Sanskrit}{-\textit{ya}-}).\footnote{This is of course an oversimplification. See \citet{Yakubovich2014}, with references, for a more nuanced discussion of the Vedic facts.} The Hittite system is more complex and includes no less than four more suffixes in addition to -\textit{iya}-\iw{Hitt}{-\textit{iye/a}-} (and \iwi{Hitt}{-\textit{āi}-} < \iwi{PIE}{*-\textit{ah\textsubscript{2}-i̯e/o}-}): \iwi{Hitt}{-\textit{aḫḫ}-}, \iwi{Hitt}{-\textit{e}-}, \iwi{Hitt}{-\textit{ešš}-}, and \iwi{Hitt}{-\textit{nu}-} (\citealt[175--179]{Hoffner2008}). The minor Anatolian languages essentially confirm this picture, with some added complications (cf.\ \citealt[547--552]{Sasseville2021}). Most Core Indo-European branches offer little or no unambiguous evidence supporting the antiquity of these formations (except for Hitt.\ \iwi{Hitt}{-\textit{e}-} < PIE \iwi{PIE}{*-\textit{eh\textsubscript{1}-i̯e/o}-}, as generally acknowledged since \citealt{Watkins1971}). The major exceptions appear to be Greek and, precisely, Balto-Slavic. Put otherwise, the interest of this branch for the reconstruction of PIE denominatives\is{denominative!verb} can hardly be overstated.

Finally, a note on dialectology. Little needs to be said about inherited formations (*-\textit{\underline{ā}-i̯e/a}-\iw{PBSl}{*-\textit{ā-tī}, *-\textit{ā-i̯e/a}-}, *-\textit{\underline{ē}-i̯e/a}-,\iw{PBSl}{*-\textit{ē-tī}, *-\textit{ē-i̯e/a}-} *-\textit{ī}-,\iw{PBSl}{*-\textit{ī}- (non-acute)} ‘\textit{nu}-verbs’). Deverbatives\is{deverbal/deverbative!verb} in *-\textit{eh\textsubscript{1}(-i̯e/o)}-,\iw{PIE}{*-\textit{eh\textsubscript{1}-i̯e/o}-} as noted above (\ref{sec:BSleje}), display a highly specific profile in the North-Western area. It remains to be seen whether the concept of ‘North-Western Indo-European’, traditionally limited to studies of vocabulary, may be extended to morphology as well. A clear \isi{Northern Indo-European} innovation is the \isi{anticausative}-\isi{inchoative} type *\textit{bʰu\nobreakdash-n\nobreakdash-dʰ\nobreakdash-é-ti}, including its extension to deadjective\is{deadjectival!verb}  \isi{fientive}s. In this area deverbative\is{deverbal/deverbative!verb} iteratives\is{iterative} in \iwi{PIE}{*-\textit{ah\textsubscript{2}-i̯e/o}-} display a \isi{productivity} unmatched elsewhere in the family. The issue clearly deserves more work. Probable or certain Balto-Slavic innovations include the suffixes *-\textit{\underline{au̯}(i̯e/a)}-\iw{PBSl}{*-\textit{au̯-tī}, *-\textit{au̯-i̯e/a}-} and *-\textit{\underline{ō}(i̯e/a)}-\iw{PBSl}{*-\textit{ō-tī}, *-\textit{ō-i̯e/a}-} and the specialization of the \textit{ā}-presents (< post-PIE \iwi{PIE}{*-\textit{ah\textsubscript{2}-e/o}-}) as deverbatives.\is{deverbal/deverbative!verb}

\subsection{PIE background 2: Formations lost in Balto-Slavic}
\label{sec:BSlPIElost}

In my view, no more (or less) denominative\is{denominative!verb} formations than those surveyed in \ref{sec:BSldenom} need to be reconstructed for Balto-Slavic. Slavic and, especially, Baltic have a wealth of extended suffixes like Lith.\ \iwi{Lith}{-\textit{d\dotacute{e}ti}}, \iwi{Lith}{-\textit{dýti}}, Sl.\ \iwi{OCS}{-\textit{kati}} etc.; see below in Section \ref{sec:BSlafter}. They usually look secondary and, as far as I can see, there is no compelling reason to project them back into Balto-Slavic.

At least one verb appears to entail \isi{conversion}: Sl.\ \iwi{PSl}{*\textit{p\H{e}lva}} ‘chaff’ $\rightarrow$ Sl.\ *\textit{p\H{e}lti}\iw{PSl}{*\textit{p\H{e}lti}, *\textit{pȇlvǫ}} (< *\textit{p\H{e}lv-ti}; \citealt[647--648]{Koch1990}), *\textit{pȇlvǫ} ‘weed’ (OCS \iwi{OCS}{\textit{plěti}, \textit{plěvǫ}}). A \isi{conversion} analysis is, in my view, a distinct (and hitherto unexplored) possibility for other etymologically isolated primary verbs,\is{derivation!deradical} e.g., OCS \iwi{OCS}{\textit{prěti}, \textit{perǫ}} ‘fly’ (cf.\ Sl.\ \iwi{PSl}{*\textit{pero̍}} ‘feather’), \iwi{OCS}{\textit{zobati}, \textit{zobljǫ}} ‘peck’ (cf.\ Sl.\ \iwi{PSl}{*\textit{zȍbъ}, -\textit{ь}} ‘crop’). At any rate, \isi{conversion} was purely sporadic in Balto-Slavic.

From a PIE perspective, the most interesting absence is that of some \textit{i̯e/o}-subtypes.\is{*-\textit{i̯e/o}-present}  The suffix \iwi{PIE}{*-\textit{i̯e/o}-} is of course ubiquitous in Balto-Slavic denominatives,\is{denominative!verb} but always as part of longer suffixes (\iwi{PIE}{*-\textit{ah\textsubscript{2}-i̯e/o}-} etc.). What is absent are \textit{i̯e/o}-denominatives\is{denominative!suffix} \is{*-\textit{i̯e/o}-present} from consonant, *-\textit{i}-, and *-\textit{u}-stems (types Ved.\ \iwi{Sanskrit}{\textit{ápas}-} ‘work’ $\rightarrow$ \textit{apas-yáti}\iw{Sanskrit}{\textit{apasyáti}} ‘works’, \iwi{Sanskrit}{\textit{jáni}-} ‘woman, wife’ $\rightarrow$ \textit{janī-yáti}\iw{Sanskrit}{\textit{janīyáti}} ‘desires a wife’, \iwi{Sanskrit}{\textit{śatru}-} ‘enemy’ $\rightarrow$ \textit{\textit{śatrū-yáti}} \iw{Sanskrit}{\textit{śatrūyáti}} ‘is hostile’). This is not by itself noteworthy (denominative suffixes\is{denominative!suffix} may easily be lost), but each “missing type” calls for some comments.

Let us begin with consonant stems. Baltic does have denominative \textit{ia}-presents,\is{denominative!verb} e.g., \iwi{Lith}{\textit{juõkas}} ‘laughter, joke’ $\rightarrow$ \iwi{Lith}{\textit{juõktis}, -\textit{iasi}} ‘laugh’, \iwi{Lith}{\textit{šveñtas}} ‘holy’ $\rightarrow$ \iwi{Lith}{\textit{švę̃sti}, \textit{šveñčia}} ‘celebrate’. These, however, almost certainly are regularizations of the obsolete type \textit{tarýti}, -\textit{ia}\iw{OLith}{\textit{tarýti, tãria, tãr\.e}} and continue, accordingly, PIE *-\textit{ei̯e/o}-\iw{PIE}{*-\textit{éi̯e/o}-} (\ref{sec:BSli}). Note, incidentally, that this implies that Baltic has no denominatives\is{denominative!verb} of the type Gk.\ \iwi{AG}{ἄγγελος} ‘messenger’ $\rightarrow$ \iwi{AG}{ἀγγέλλω} ‘bear a message’ (< *°\textit{el-i̯e/o}-), with deletion of the thematic vowel. Slavic has a large class of “expressive” primary verbs\is{derivation!deradical}  such as OCS \iwi{OCS}{\textit{glagolati}, \textit{glagolje}-} ‘speak’. Some of them could be interpreted as denominatives\is{denominative!verb} (cf.\ \iwi{OCS}{\textit{glagolъ}} ‘speech’), but most certainly not, cf.\ \citet[328--346]{Vaillant1966}. Note, finally, that Balto-Slavic has no historically complex suffixes like Gk.\ -ιζω \iw{AG}{-ιζω, -ιζειν} (< *-\textit{id-i̯e/o}-) or Toch.\ \iwi{CToch}{*-\textit{ññä/æ}-} (< *-\textit{n-i̯e/o}-).

The lack of a special suffix for \textit{i}-stem denominatives\is{denominative!verb}  is fully expected. The development PIE *-\textit{ei̯e/o}-\iw{PIE}{*-\textit{éi̯e/o}-} > Bl.-Sl.\ *-\textit{ī}-\iw{PBSl}{*-\textit{ī}- (non-acute)} must have started with an early raising \mbox{*-\textit{ei̯e/o}-} > \mbox{-\textit{ii̯e/o}-,} which entailed merger with \textit{i}-stem denominatives\is{denominative!verb}  in \iwi{PIE}{*-\textit{i-i̯e/o}-} (see above \tabref{tab:ejeoBSl}). This is confirmed by Slavic, where \textit{i}-verbs build denominatives\is{denominative!verb} from -\textit{o}- and -\textit{i}-stems alike. The Baltic type \textit{dalýti}, -\textit{ìja}\iw{Lith}{\textit{dalýti, -ìja, -ìjo}} does not go back to \textit{i}-stem denominatives\is{denominative!verb} (cf.\ \citealt{VillanuevaSvensson2023b}).

The \textit{u}-stems are more interesting. We do not only lack reflexes of \iwi{PIE}{*-\textit{u-i̯e/o}-}; the stem vowel -\textit{u}-\iw{PBSl}{*-\textit{u}-} is absent in denominatives:\is{denominative!verb} Lith.\ \textit{saldùs}\iw{Lith}{\textit{saldùs} m., \textit{saldì} f.} ‘sweet’ (OCS \iwi{OCS}{\textit{sladъkъ}}) $\rightarrow$ \textit{sald-\dotacute{e}ti}\iw{Lith}{\textit{sald\dotacute{e}ti}} ‘be, become sweet’, \textit{sáld-inti}\iw{Lith}{\textit{sáldinti}} ‘sweeten’. In recent years the matter has received two different interpretations. \citet[117--118]{Ostrowski2006} suggests that we are dealing with Caland system relics;\is{Caland!system (CS)}  see also \citet[261--262, and \textit{passim}]{Majer2021} for discussion. \citet[6--13]{Hill2012}, on the other hand, has argued for a \isi{Northern Indo-European} rule of \textit{u}-deletion before \textit{jod}. At this point it should be noted that \iwi{PIE}{*-\textit{u}-} is also absent from the feminine of \textit{u}-stem adjectives: Lith.\ \textit{saldì}, -\textit{džios}\iw{Lith}{\textit{saldùs} m., \textit{saldì} f.} (contrast Ved.\ \textit{ur-ú-ḥ}\iw{Sanskrit}{\textit{urú}-} : \textit{ur-v-ī́},\iw{Sanskrit}{\textit{urvī́}-} Gk.\ \iwi{AG}{γλυκύς} : \iwi{AG}{γλυκει̃α} < -\textit{eu̯-i̯a}). As a third option I would suggest analogy with \textit{o}-stem adjectives, where the thematic vowel is (synchronically) deleted: Lith.\ \iwi{Lith}{\textit{plónas}} ‘thin’ $\rightarrow$ \textit{plon-\dotacute{e}ti}\iw{Lith}{\textit{plon\dotacute{e}ti}} ‘become thin’, \textit{plón-inti}\iw{Lith}{\textit{plóninti}} ‘make thin’. The issue clearly deserves more work.

\subsection{Balto-Slavic denominatives from a typological perspective}
\label{sec:BSltypolog}

In this section I will briefly survey what Balto-Slavic tells us about the way denominatives \is{denominative!verb} arise and develop. The inference of what we have seen seems to be that (almost) anything goes. Denominatives\is{denominative!verb} may be inherited (e.g., *-\textit{ei̯e/o}-,\iw{PIE}{*-\textit{éi̯e/o}-} \iwi{PIE}{*-\textit{eh\textsubscript{1}-i̯e/o}-}), created anew (*-\textit{\underline{au̯}-i̯e/a}-),\iw{PBSl}{*-\textit{au̯-tī}, *-\textit{au̯-i̯e/a}-} or transformed in different ways (e.g., Lith.\ \iwi{Lith}{-\textit{inti}, -\textit{ina}} and \iwi{Lith}{-\textit{ýti}, -\textit{ìja}}). The Balto-Slavic picture is perhaps more interesting from the viewpoint of semantics and, especially, the interplay of denominatives\is{denominative!verb} and deverbatives.\is{deverbal/deverbative!verb}

I am not aware of detailed studies on denominative semantics from a historical perspective. PIE *-\textit{ei̯e/o}-\iw{PIE}{*-\textit{éi̯e/o}-} (Bl.-Sl.\ *-\textit{ī}-\iw{PBSl}{*-\textit{ī}- (non-acute)}) was often (though not exclusively) \isi{transitive}, but the semantics of *-\textit{\underline{ā}-i̯e/a}-,\iw{PBSl}{*-\textit{ā-tī}, *-\textit{ā-i̯e/a}-} *-\textit{\underline{ō}-i̯e/a}-,\iw{PBSl}{*-\textit{ō-tī}, *-\textit{ō-i̯e/a}-} and *-\textit{\underline{au̯}-i̯e/a}-\iw{PBSl}{*-\textit{au̯-tī}, *-\textit{au̯-i̯e/a}-} would definitely benefit from more work. Deadjectives \is{deadjectival!verb} in *-\textit{\underline{ē}-i̯e/a}-\iw{PBSl}{*-\textit{ē-tī}, *-\textit{ē-i̯e/a}-} developed from \isi{essive}s to \isi{fientive}s in both Baltic and Slavic, but the interplay between both values is well known. As for the interplay between denominatives\is{denominative!verb} and deverbatives,\is{deverbal/deverbative!verb} it has been rampant. This may be a consequence, in part, of the architecture of the Balto-Slavic verb, which can be seen as a network of interrelated verbs, each of them filling the slot of a given Aktionsart. See \tabref{tab:Lithverbs}.

\begin{table}
\caption{Lithuanian verb system}
\label{tab:Lithverbs}
    \begin{tabularx}{\textwidth}{llllll}
    \lsptoprule
    & I & II & III & IV & V \\
    & Unmarked & Stative-\is{stative} & \is{inchoative}Inchoative- & Causative\is{causative} & Iterative\is{iterative} \\
    & (\isi{transitive})  &\isi{durative} & \isi{anticausative} & & (etc.)\\
    \midrule
    & \textit{a}-, \textit{ia}- & 2\textsuperscript{nd} stem & -\textit{n}-, \textit{sta}- & -\textit{ýti}, -\textit{ìnti} & -\textit{ýti}, -\textit{óti}, \\
    &presents &*-\textit{ē}-, *-\textit{ā}- &presents & &\textit{in\dotacute{e}ti}, etc. \\
   \midrule
   
    strew & \iwi{Lith}{\textit{bérti, -ia}} & \iwi{Lith}{\textit{byr\dotacute{e}ti, bỹra}} & \iwi{Lith}{\textit{bìrti, bỹra}} & \iwi{Lith}{\textit{bìrinti}} & \iwi{Lith}{\textit{barstýti}},\\
    (\iwi{PIE}{*\textit{bʰer}-}) &&&\textit{(< *biñra)}&&\iwi{Lith}{\textit{birénti}},\\
    &&&&&\iwi{Lith}{\textit{birsnóti}}\\
    
    lift & \iwi{Lith}{\textit{kélti, kẽlia}} & \textit{pa-kyl\dotacute{e}ti}\iw{Lith}{\textit{pakyl\dotacute{e}ti}} & \iwi{Lith}{\textit{kìlti, kỹla}} & \iwi{Lith}{\textit{kéldinti}},& \iwi{Lith}{\textit{kilnóti}},\\
    (\iwi{PIE}{*\textit{kelh\textsubscript{3}}-}) &&&\textit{(< *kiñla)}&\iwi{Lith}{\textit{kìldyti}},&\iwi{Lith}{\textit{kilóti}},\\
    &&&&\iwi{Lith}{\textit{kìldinti}} & \iwi{Lith}{\textit{kilsúoti}},\\
    &&&&&\iwi{Lith}{\textit{kilstel\dotacute{e}ti}}\\
     
    peck & \iwi{Lith}{\textit{lèsti, lẽsa}} &  &  & \iwi{Lith}{\textit{lèsinti}} & \iwi{Lith}{\textit{lesióti}}\\
    (\iwi{PIE}{*\textit{les}-})&&&&&\\
    
    sleep  &  & \iwi{Lith}{\textit{miegóti, -a}} & \textit{už-mìgti},\iw{Lith}{\textit{užmìgti, -miñga}} & \iwi{Lith}{\textit{migdýti}} & \\
    (\iwi{PIE}{*\textit{mei̯gʰ}-}) &&&\textit{-miñga}&&\\
    \lspbottomrule
    \end{tabularx}
\end{table}

It is also important to note that Balto-Slavic, unlike other branches, has preserved the autonomy of the root up to very recently (in Baltic to some degree up to this day). At any rate, it is evident that such a system provided abundant grounds for mutual contamination between denominatives\is{denominative!verb} and deverbatives.\is{deverbal/deverbative!verb} Interestingly, not all theoretical configurations took place with the same intensity:

\begin{enumerate}
    \item “Denominative\is{denominative!verb} $\rightarrow$ secondary verbal formation\is{derivation!stem-based} (\isi{iterative}, \isi{causative}, etc.)” is the most common development. It happened in the case of PIE \iwi{PIE}{*-\textit{ah\textsubscript{2}-i̯e/o}-}, \iwi{PIE}{*-\textit{eh\textsubscript{1}-i̯e/o}-}, Bl.-Sl.\ *-\textit{\underline{ō}-i̯e/a}-,\iw{PBSl}{*-\textit{ō-tī}, *-\textit{ō-i̯e/a}-} *-\textit{\underline{au̯}-i̯e/a}-,\iw{PBSl}{*-\textit{au̯-tī}, *-\textit{au̯-i̯e/a}-} and Bl.\ *-\textit{\underline{in}ti}, *-\textit{ina},\iw{PB}{*-\textit{inti}, *-\textit{ina}} if at very different dates. If my account is correct, in the case of the Balto-Slavic \textit{ā}-presents and the Slavic type \textit{rinǫti}, \textit{rine}-\iw{OCS}{\textit{rinǫti, rine-, rinǫxъ}, ppp \textit{rinovenъ}} an original denominative\is{denominative!verb} type survives as deverbative\is{deverbal/deverbative!verb} alone. Another generalization is that once deverbatives\is{deverbal/deverbative!verb} are established in the system, they easily evolve towards more complicated morphology, whereas denominatives\is{denominative!verb} tend to keep simpler morphology.

    \item “Denominative\is{denominative!verb} $\rightarrow$ (unmarked) primary verb”\is{derivation!deradical} certainly occurs in isolated verbs. What seems to be genuinely rare is the creation of primary verbal\is{derivation!deradical} classes out of denominatives.\is{denominative!verb} Potential cases (primary \textit{ē}-statives,\is{stative} \textit{ā}-aorist) almost certainly passed through a stage involving secondary verbal\is{derivation!verbal} derivation.\is{derivation!stem-based}

    \item “Secondary verbal formation\is{derivation!stem-based} $\rightarrow$ denominative”\is{denominative!verb} appears to be very rare. The only certain case is the \isi{Northern Indo-European} \isi{fientive}  type *\textit{gʰru\nobreakdash-m-bʰ-e/o}-, which probably responded to a highly specific situation. Slavic \textit{i}-\isi{factitive}s are easily explained as an extension of \isi{transitive} desubstantives, with deverbal\is{deverbal/deverbative!verb} \isi{causative}s playing at best a supporting role. If my account is correct \isi{causative}s like Lith.\ \textit{dẽg-inti}, -\textit{ina}\iw{Lith}{\textit{dẽginti, -ina, -ino}} are a secondary extension of \isi{factitive} deadjectives.\is{deadjectival!verb}

    \item Finally, “(unmarked) primary\is{derivation!deradical} verb $\rightarrow$ denominative”\is{denominative!verb} is unattested. It is noteworthy that secondary\is{derivation!stem-based} deverbal\is{deverbal/deverbative!verb} formations appear to go back to denominatives\is{denominative!verb} much more frequently than to \isi{reanalysis} of primary verb morphology.\is{derivation!deradical} 
\end{enumerate}

If these observations are accurate, they are potentially interesting for work on other branches, especially in the case of formations that, in theory, could be of both denominative\is{denominative!verb} and deverbative\is{deverbal/deverbative!verb} origin (e.g., \isi{factitive}-\isi{causative}s). It is evident, however, that they should be tested against a more extensive and clearly delimited dataset before they can be regarded as general tendencies of language change.

\subsection{After Balto-Slavic}
\label{sec:BSlafter}

The Balto-Slavic verbal system sketched above (\ref{sec:BSltypolog}) provided abundant grounds for the creation of new complex suffixes, semantic specialization, etc. Also, of course, for contamination with denominatives.\is{denominative!verb} Here we are only interested in the development of the system after the breakup of Balto-Slavic, with some notes on possible areal phenomena:

\begin{enumerate}
    \item The system was not only preserved, but actually reinforced in Baltic (note especially, but not only the development of the \textit{ī}-verbs). In addition, Baltic is rich in semantically and formally specialized minor subtypes (“\isi{intensive}s”, “diminutives”, etc.). It is interesting to note that on this score Baltic contrasts with the major Indo-European branches of Europe. From an areal perspective, it is tempting to attribute the divergent development of Baltic, in part, to Uralic influence (as suggested by \citealt{Pakerys2018} for the \isi{causative}).

    \item The system was eliminated in Slavic. Types like the \textit{nǫ}-\isi{inchoative}s, the \textit{ě}-statives,\is{stative} or the \textit{i}-causatives\is{causative} are still abundantly represented, but are not productive\is{productivity} anymore. The \textit{a}-iteratives,\is{iterative} of course, became immensely productive,\is{productivity} but as part of a new system of grammatical aspect. From a typological point of view late Slavic came to be aligned with the modern Germanic and Romance languages in disfavoring secondary\is{derivation!stem-based}  verbal derivation\is{derivation!verbal} altogether. The Slavic development may indeed have been scattered by contacts with these languages.
\end{enumerate}

If these observations are correct, the recent development of Baltic and Slavic awaits further work. The matter, in addition, has an obvious interest for the areal dimension of language change. Be it as it may, the divergent tendencies towards complication in Baltic and simplification in Slavic are clearly seen in the denominative\is{denominative!verb} system as well.

To conclude, I hope to have shown that Baltic and Slavic denominatives\is{denominative!verb} are not only interesting for their own sake; they have a genuine contribution to make to the reconstruction of the PIE system.

\section*{Acknowledgements}

The final version of this paper has benefited from the comments of Marek Majer, Jurgis Pakerys, Viktoria Reiter and an anonymous reviewer of an earlier version of it. My thanks must be extended to the participants in the Workshop and to the organizers. Needless to say, I alone am responsible for the views presented here. 
 
 
\section*{Abbreviations}
\begin{tabularx}{.48\textwidth}{lQ}
Aor. & Aorist\\
Bl.\ & Baltic\\
Bl.-Sl.\ & Balto-Slavic\\
CSl.\ & Common Slavic\\
Fem. & Feminine\\
Gen. & Genitive\\
Gk.\ & Greek\\
Gmc. & Germanic\\
Go. & Gothic\\
Hitt. & Hittite\\
Impf. & Imperfective\\
Inf. & Infinitive\\
Intr. & Intransitive\\
Iter. & Iterative\\
Lat.\ & Latin\\
Latv.\ & Latvian\\
Lith.\ & Lithuanian\\
OCS & Old Church Slavonic\\
\end{tabularx}
\begin{tabularx}{.48\textwidth}{lQ}
OIr.\ & Old Irish\\
OLith.\ & Old Lithuanian\\
ON & Old Norse\\
OPr.\ & Old Prussian\\
Pass. & Passive\\
PIE & Proto-Indo-European\\
Pl.\ & Plural\\
Pl.t. & Pluralia tantum\\
Pol. & Polish\\
Pret. & Preterite\\
Ptcp. & Participle\\
Refl. & Reflexive\\
Sg.\ & Singular\\
Sl.\ & Slavic\\
Toch. & Tocharian\\
Tr. & Transitive\\
Ved.\ & Vedic
\end{tabularx}



{\sloppy\printbibliography[heading=subbibliography,notkeyword=this]}
\end{document}
